\documentclass[11pt,a4paper]{article}

% ----------- Packages -----------
\usepackage[utf8]{inputenc}
\usepackage[T1]{fontenc}
\usepackage[english]{babel}
\usepackage{lmodern}
\usepackage{geometry}
\geometry{margin=2.5cm}
\usepackage{setspace}
\onehalfspacing

\usepackage{hyperref}
\hypersetup{
  colorlinks=true,
  linkcolor=black,
  urlcolor=blue,
  citecolor=black,
  pdfauthor={},
  pdftitle={Radio-AI Product Whitepaper},
  pdfsubject={},
  pdfkeywords={}
}

\usepackage{graphicx}
\usepackage{caption}
\usepackage{subcaption}
\usepackage{amsmath,amssymb}
\usepackage{enumitem}
\usepackage{titlesec}
\usepackage{xcolor}
\usepackage{listings}
\usepackage{booktabs}
\usepackage{float}

% ---- Styling sections (optional, clean tech look) ----
\titleformat{\section}{\large\bfseries}{\thesection}{0.7em}{}
\titleformat{\subsection}{\normalsize\bfseries}{\thesubsection}{0.6em}{}
\titleformat{\subsubsection}{\normalsize\itshape}{\thesubsubsection}{0.5em}{}

% ----------- Title -----------
\title{\textbf{Radio-AI Product Whitepaper}\\[4pt]
\large Personalized AI-Generated Radio Shows with Thematic Advertising}
\author{%
  Samuele Tamborrino \\
  Marco Ravaioli
}
\date{\today}

\begin{document}

\maketitle
\thispagestyle{empty}
\vfill
\begin{center}
    \textit{Version 0.1 -- Draft for internal review}
\end{center}
\newpage

\tableofcontents
\newpage

% ---------------------------------
\section{Introduzione \& Obiettivi}

\subsection{Contesto}

L’ecosistema audio contemporaneo è dominato da tre formati principali: lo streaming musicale on-demand, i podcast e la radio tradizionale.  
Le piattaforme musicali offrono cataloghi pressoché infiniti, ma si limitano a suggerire brani tramite algoritmi: non costruiscono narrazioni, non creano atmosfera e non offrono contesto personalizzato.  
I podcast rappresentano contenuti strutturati e di alta qualità, ma sono statici, non reagiscono all’utente e richiedono una produzione significativa.  
La radio tradizionale mantiene il valore della spontaneità e della compagnia, ma non può adattarsi ai gusti, al momento o alle emozioni di chi ascolta.

L’avanzamento recente di Large Language Models (LLM), Text-to-Speech neurale (TTS) e sistemi di streaming a bassa latenza abilita un nuovo paradigma: contenuti audio generati in tempo reale, personalizzati per ogni singolo utente.  
In questo contesto nasce \textbf{Radio-AI}, una piattaforma che fonde musica reale, narrazione dinamica e intelligenza artificiale in un’unica esperienza radiofonica personalizzata.

\subsection{Problem Statement}

Nonostante la grande disponibilità di contenuti audio, esistono ancora importanti limiti per ascoltatori, creatori e brand.

\paragraph{Ascoltatori}
\begin{itemize}
    \item Mancanza di \textbf{personalizzazione contestuale}: la musica è personalizzata, la narrazione no.
    \item Assenza di \textbf{compagnia adattiva}: la radio tradizionale ha una voce umana, ma non può adattarsi all’utente o alla sua situazione.
    \item Difficoltà ad integrare musica, informazioni culturali e intrattenimento in un unico flusso coerente.
    \item Incapacità delle piattaforme attuali di offrire un’esperienza che combini musica, curiosità verificate e momenti narrativi personalizzati.
\end{itemize}

\paragraph{Creator}
\begin{itemize}
    \item La produzione di contenuti audio richiede tempo, risorse e capacità tecniche.
    \item Limitata possibilità di creare contenuti personalizzati su larga scala.
    \item Difficoltà nel raggiungere nicchie specifiche con contenuti realmente su misura.
\end{itemize}

\paragraph{Brand}
\begin{itemize}
    \item Le attuali soluzioni di advertising audio si basano su \textbf{target demografici} poco contestuali.
    \item Mancanza di strumenti per inserire messaggi pubblicitari in \textbf{momenti narrativi coerenti}.
    \item Crescente necessità di forme pubblicitarie \textbf{native, tematiche e orientate all’emozione} anziché al profilo utente.
\end{itemize}

Radio-AI affronta queste criticità offrendo un sistema in grado di generare dirette personalizzate, integrando contenuti verificati, narrazione flessibile e inserzioni tematiche coerenti con l’esperienza dell’ascolto.

\subsection{Obiettivi del Prodotto}

Radio-AI nasce con obiettivi chiari, suddivisi in tre macro-aree: esperienza utente, architettura tecnica e modello di business.

\paragraph{Obiettivi di User Experience}
\begin{itemize}
    \item Produrre una \textbf{diretta radiofonica personalizzata} a partire da un unico prompt e/o da una playlist selezionata.
    \item Offrire una sensazione di \textbf{presenza e compagnia} tramite uno speaker AI con voce e tono personalizzabili.
    \item Integrare in modo fluido \textbf{curiosità e informazioni reali} all’interno dell’intrattenimento musicale.
    \item Permettere \textbf{modifiche in tempo reale} del prompt con aggiornamento immediato della diretta.
    \item Abilitare funzionalità social come \textbf{Radio Rooms}, Quote Cards condivisibili e profili pubblici.
\end{itemize}

\paragraph{Obiettivi Tecnici}
\begin{itemize}
    \item Progettare un’architettura scalabile e modulare basata su generazione audio \textbf{chunked} a bassa latenza.
    \item Garantire la \textbf{sincronizzazione} tra voce AI e musica riprodotta tramite servizi esterni (Spotify / Apple Music).
    \item Implementare un \textbf{motore di fact-checking} basato su pipeline RAG per contenuti verificati e affidabili.
    \item Supportare interazioni distribuite e in tempo reale con un sistema WebSocket per le \textbf{Radio Rooms}.
    \item Assicurare affidabilità, resilienza e costi sostenibili attraverso caching, autoscaling e CDN.
\end{itemize}

\paragraph{Obiettivi di Business}
\begin{itemize}
    \item Creare un modello freemium con funzionalità premium (voci avanzate, dirette più lunghe, qualità audio superiore).
    \item Lanciare un \textbf{Radio Ads Manager} capace di targettizzare temi, atmosfere ed emozioni anziché dati demografici.
    \item Costruire una piattaforma social di contenuti generati dagli utenti, favorendo \textbf{engagement e retention}.
    \item Sostenere la crescita tramite una combinazione di revenue B2C (abbonamenti) e B2B (advertising contestuale).
\end{itemize}% ---------------------------------
\section{System Overview}

\subsection{Sintesi dell'Esperienza Utente}

Radio-AI offre un flusso d’uso semplice e unificato basato su un’unica barra di prompt, capace di comprendere preferenze, tono desiderato, temi, rubrica, mood e indicazioni musicali. Il sistema guida l’utente in tre passaggi principali:

\begin{enumerate}
    \item \textbf{Input dell’utente}:  
    L’utente può inserire un prompt libero (breve o articolato) e, se desidera, selezionare una playlist da servizi esterni come Spotify o Apple Music. Se non viene scelta una playlist, l’AI ne genera una coerente con il tema richiesto.

    \item \textbf{Generazione della diretta}:  
    Il Prompt Interpreter analizza l’intento, estrae le categorie semantiche (tema, rubriche, curiosità reali da includere, mood, stile dello speaker, preferenze musicali), produce una scaletta coerente e genera uno script suddiviso in \emph{chunk} narrativi.  
    Parallelamente, la musica viene gestita tramite il servizio dell’utente, mentre il motore TTS produce la voce dello speaker e l’orchestratore audio sincronizza parlato e brani. Nella modalità senza playlist pre-esistente, prima dell’avvio definitivo l’utente viene sempre invitato ad aggiungere eventuali brani “imprescindibili”, che vengono integrati nella scaletta generata dall’AI.

    \item \textbf{Riproduzione e interazione}:  
    L’utente ascolta una diretta radiofonica personalizzata, con possibilità di mettere in pausa, modificare il prompt in tempo reale, partecipare o creare Radio Rooms, salvare la puntata, generare Quote Cards e condividere l’esperienza.
\end{enumerate}

Questo approccio combina la spontaneità della radio, la personalizzazione dello streaming moderno e la flessibilità narrativa del contenuto generato da AI.
Inoltre, l’esperienza utente si articola in due modalità principali: Modalità A, in cui l’utente seleziona una playlist esistente (Spotify/Apple Music) e la puntata viene costruita usando solo quella selezione, e Modalità B, in cui la playlist viene generata dall’AI a partire dal prompt, con la possibilità di aggiungere manualmente brani specifici prima dell’avvio. In entrambe le modalità, l’utente può controllare il livello di presenza dello speaker AI, fino a configurazioni quasi “solo musica” o interamente musicali con brevi citazioni o intermezzi minimali.

\subsection{Funzionalità di Alto Livello}

Radio-AI integra un insieme di funzionalità trasversali che costituiscono la piattaforma:

\begin{itemize}
    \item \textbf{AI Speaker Personalizzabile(anche disattivabile)}: voce sintetica con timbro, tono e stile configurabili; possibilità di usare voci premium ispirate a personaggi famosi (previa licenza) e di ridurre o azzerare del tutto la presenza dello speaker, per supportare puntate “solo musica” o basate principalmente su citazioni e stacchi minimali.

    \item \textbf{Curiosità e Contenuti Verificati}:  
    Motore di fact-checking che valida informazioni culturali, storiche o musicali tramite pipeline RAG e fonti affidabili.

    \item \textbf{Generazione Playlist}:  
    Modalità AI-driven per selezione musicale coerente con tema e mood, con possibilità di inserire brani specifici.

    \item \textbf{Script Engine a Chunk}:  
    Generazione incrementale di parlato e rubriche per garantire latenza ridotta e modifiche live.

    \item \textbf{Radio Rooms}:  
    Ascolto condiviso in tempo reale tramite infrastruttura WebSocket, con audio sincronizzato tra più dispositivi.

    \item \textbf{Quote Cards}:  
    Estratti testuali generati dalla trascrizione della diretta, trasformabili in card grafiche personalizzabili e condivisibili.

    \item \textbf{Profilo Utente e Cronologia}:  
    Archivio delle puntate generate, prompt salvati, Quote Cards, preferenze e contenuti condivisi.

    \item \textbf{Ads Manager}:  
    Sistema di advertising tematico che inserisce spot coerenti con atmosfera, tema e contesto narrativo della puntata.

    \item \textbf{Editing in Tempo Reale}:  
    L’utente può modificare prompt, tono, rubriche o musica mentre la puntata è in corso; il sistema rigenera i chunk successivi mantenendo la continuità narrativa.
\end{itemize}

\subsection{Ambito di Questo Documento}

Questo whitepaper descrive la progettazione tecnica, l’architettura software e le considerazioni ingegneristiche alla base di Radio-AI.  
Nello specifico, il documento copre:

\begin{itemize}
    \item la visione di alto livello del sistema e dei flussi utente;
    \item l’architettura generale e i moduli software principali;
    \item i componenti AI (LLM, TTS, fact-checking, rendering audio);
    \item i flussi logici delle funzionalità core (generazione puntata, editing live, Radio Rooms, Ads Manager);
    \item lo stack tecnologico raccomandato;
    \item i requisiti infrastrutturali e le strategie di scalabilità;
    \item il modello API di alto livello;
    \item la roadmap tecnica di sviluppo.
\end{itemize}

Non rientrano nell’ambito di questo documento:

\begin{itemize}
    \item specifiche legali dettagliate su licenze musicali e utilizzo di voci di personaggi famosi;
    \item analisi economico-finanziarie approfondite (addressable market, competitor analysis);
    \item design grafico dell’applicazione mobile o web al livello UI pixel-perfect;
    \item implementazione completa o codice sorgente dei microservizi.
\end{itemize}

L’obiettivo è fornire una base chiara e strutturata per ingegneri, CTO, partner tecnici e potenziali collaboratori, facilitando una valutazione realistica della fattibilità, della scalabilità e delle componenti tecnologiche necessarie.
% ---------------------------------
\section{Architecture Overview}

\subsection{Global Architecture Diagram}

La Figura~\ref{fig:architecture} illustra la visione d’insieme dell’architettura Radio-AI.  
Il sistema è composto da tre macro–aree principali:

\begin{itemize}
    \item \textbf{User Input Layer} — il livello in cui l’utente definisce prompt, preferenze musicali e impostazioni vocali.
    \item \textbf{AI Pipeline Layer} — il cuore logico del sistema, responsabile dell’interpretazione del prompt, della generazione dello script, del fact-checking e della produzione audio.
    \item \textbf{Output Layer} — il livello di playback che combina musica esterna e contenuti vocali generati dall’AI, includendo anche funzionalità di sincronizzazione e condivisione.
\end{itemize}

La pipeline è progettata in modalità \textbf{loosely coupled micro-services}, consentendo scalabilità orizzontale, caching dei risultati e aggiornamenti indipendenti dei moduli principali.

\begin{figure}[H]
    \centering
    \includegraphics[width=0.9\textwidth]{fig/architecture-radio-ai.png}
    \caption{Radio-AI: architettura globale dalla fase di input fino alla generazione dell’esperienza audio finale.}
    \label{fig:architecture}
\end{figure}

\subsection{Main Components}

L’architettura Radio-AI si compone dei seguenti blocchi principali.

\paragraph{1. API Gateway}
Punto di ingresso unificato per l’app client (mobile o web). Gestisce:
\begin{itemize}
    \item autenticazione e autorizzazione;
    \item rate limiting;
    \item routing verso i microservizi interni.
\end{itemize}

\paragraph{2. Prompt Interpreter}
Modulo basato su LLM che elabora il prompt naturale dell’utente.  
Estrae intenzioni preliminari e costruisce un JSON strutturato contenente:
\begin{itemize}
    \item tema della puntata;
    \item mood e stile dello speaker;
    \item rubriche richieste;
    \item preferenze musicali o vincoli;
    \item richieste di contenuto culturale o informativo.
\end{itemize}

\paragraph{3. Intent Extractor}
Raffina e valida le intenzioni estratte dal prompt.  
Applica euristiche, controlli di coerenza e normalizzazione semantica per garantire input standardizzati al resto della pipeline.

\paragraph{4. Playlist Builder}
Genera una playlist coerente con il tema, oppure integra la playlist esterna selezionata dall’utente.  
Nella modalità “senza playlist”, permette l’inserimento manuale di brani prima della generazione.

\paragraph{5. Fact-Checking Engine}
Pipeline RAG (Retrieval-Augmented Generation) che verifica contenuti culturali:  
\begin{itemize}
    \item recupera informazioni da fonti affidabili (database musicali, knowledge graphs, enciclopedie);
    \item scarta o riscrive contenuti non verificabili;
    \item fornisce uno strato di «verità» sopra cui lo Script Generator può costruire la narrazione.
\end{itemize}

\paragraph{6. Script Generator}
Genera l’intero script della puntata, dividendolo in \textbf{chunk narrativi} (30–90 secondi).  
Produce:
\begin{itemize}
    \item intro e apertura;
    \item rubriche e approfondimenti;
    \item collegamenti tra i brani;
    \item curiosità verificate;
    \item eventuali elementi fiction (telefonate, sketch);
    \item outro e chiusura.
\end{itemize}

\paragraph{7. Chunked Renderer}
Coordina la generazione dei blocchi audio.  
Esegue:
\begin{itemize}
    \item conversione testo $\rightarrow$ TTS;
    \item caching e anticipazione dei chunk successivi;
    \item versioni “Fast Render” e “High Quality Render”;
    \item gestione della rigenerazione in caso di modifica del prompt.
\end{itemize}

\paragraph{8. TTS Engine}
Modulo di sintesi vocale basato su provider esterni (Azure TTS, ElevenLabs o equivalenti), con possibilità di:
\begin{itemize}
    \item personalizzare timbro, tono, postura vocale;
    \item selezionare voci premium o ispirate a personaggi famosi.
\end{itemize}

\paragraph{9. Audio Playback \& Sync Layer}
Livello che combina:
\begin{itemize}
    \item musica riprodotta via Spotify/Apple Music dal dispositivo dell’utente;
    \item voce AI proveniente dal CDN;
    \item sincronizzazione temporale tra i due flussi;
    \item gestione delle Radio Rooms via WebSocket.
\end{itemize}

\paragraph{10. Storage \& CDN}
Include:
\begin{itemize}
    \item database relazionale per utenti, episodi, campagne e cronologia;
    \item storage oggetti per chunk audio, Quote Cards e media generati;
    \item CDN per distribuzione rapida dei contenuti.
\end{itemize}

\paragraph{11. Radio Ads Manager}
Modulo separato per la gestione di campagne pubblicitarie.  
Fornisce funzionalità di:
\begin{itemize}
    \item creazione obiettivo;
    \item targeting tematico;
    \item generazione spot via AI;
    \item posizionamento ottimale dentro l’episodio;
    \item reportistica dettagliata.
\end{itemize}

\subsection{Data Flow}

Il flusso dei dati in Radio-AI segue una sequenza strutturata:

\begin{enumerate}
    \item \textbf{Input}  
    L’utente invia un prompt e, se necessario, una playlist esterna.

    \item \textbf{Interpretazione}  
    Il Prompt Interpreter analizza il testo e produce un modello strutturato delle intenzioni.

    \item \textbf{Validazione}  
    L’Intent Extractor pulisce, normalizza e conferma le istruzioni operative.

    \item \textbf{Preparazione contenuti}  
    Playlist Builder crea la selezione musicale oppure integra quella fornita.

    \item \textbf{Verifica informazioni}  
    Il Fact-Checking Engine recupera e valida dati culturali e storici richiesti.

    \item \textbf{Generazione script}  
    Lo Script Generator produce la narrazione completa, divisa in chunk.

    \item \textbf{Rendering}  
    Il Chunked Renderer crea e sincronizza i blocchi audio tramite il TTS Engine.

    \item \textbf{Playback}  
    Il dispositivo dell’utente riproduce musica e voce AI in sincronia, anche dentro Radio Rooms.

    \item \textbf{Persistenza}  
    La puntata, i chunk generati, i metadati, la trascrizione e i tag vengono salvati per riutilizzo, ricerca e condivisione.
\end{enumerate}

Questa pipeline permette un avvio rapido della diretta, una rigenerazione dinamica durante l’ascolto e scalabilità verticale e orizzontale sui moduli più intensivi (LLM, TTS e rendering).
% ---------------------------------
\section{Software Modules}

\subsection{Core Backend Services}

Il backend di Radio-AI è progettato come un insieme di microservizi indipendenti ma coordinati, ciascuno responsabile di una parte specifica della pipeline di generazione e riproduzione della puntata. Questa architettura favorisce scalabilità, resilienza, aggiornabilità indipendente e un chiaro controllo delle risorse computazionali.

Di seguito vengono descritti i principali servizi core.

\paragraph{API Gateway}
È il punto di ingresso centralizzato dell’intera piattaforma.  
Responsabilità principali:
\begin{itemize}
    \item Gestione dell’autenticazione (OAuth2 con provider esterni, JWT).
    \item Routing delle richieste verso i microservizi interni.
    \item Enforcement di rate limiting e throttling.
    \item Logging, auditing e tracciamento distribuito delle richieste.
    \item Prima validazione dei payload ricevuti dal client.
\end{itemize}

\paragraph{Session Service}
Mantiene lo stato delle sessioni di generazione e ascolto.  
Responsabilità:
\begin{itemize}
    \item Gestione dello stato della puntata in corso (chunk attuale, timestamp, modalità di ascolto).
    \item Sincronizzazione con le Radio Rooms.
    \item Gestione dei checkpoint per la rigenerazione in caso di modifiche al prompt.
    \item Associazione tra utente, episodio e dispositivo.
\end{itemize}

\paragraph{Prompt Interpreter Service}
È il primo modulo di AI nella catena logica.  
Responsabilità:
\begin{itemize}
    \item Analisi del prompt naturale fornito dall’utente.
    \item Identificazione delle categorie semantiche richieste: tema, mood, stile, rubriche, curiosità culturali, preferenze musicali.
    \item Parsing e output in formato JSON strutturato.
    \item Normalizzazione dei parametri per uso interno.
\end{itemize}

\paragraph{Playlist Service}
Gestisce la definizione musicale dell’episodio.  
Responsabilità:
\begin{itemize}
    \item Parsing e validazione di playlist provenienti da Spotify/Apple Music.
    \item Generazione autonoma di playlist in modalità “senza playlist”.
    \item Gestione dell’inserimento manuale di brani da parte dell’utente.
    \item Uniformazione dei metadati musicali alla timeline dell’episodio.
\end{itemize}

\paragraph{Script Generator Service}
Costruisce l’intera narrazione della puntata.  
Responsabilità:
\begin{itemize}
    \item Creazione dello script completo (intro, rubriche, curiosità, sketch, outro).
    \item Suddivisione in \textbf{chunk} narrativi da 30--90 secondi.
    \item Generazione coerente con tema, tono, durata e mood.
    \item Integrazione dei risultati del Fact-Checking Service.
    \item Supporto alla rigenerazione di singoli chunk dopo modifiche al prompt.
\end{itemize}

\paragraph{Fact-Checking Service}
È il modulo responsabile della verifica delle informazioni reali.  
Responsabilità:
\begin{itemize}
    \item Recupero di informazioni da knowledge graph, API musicali, enciclopedie e dataset certificati.
    \item RAG pipeline per verificare, confermare o riformulare contenuti non verificabili.
    \item Memorizzazione nella cache delle informazioni più richieste.
    \item Fornitura allo Script Generator di dati affidabili e contestuali.
\end{itemize}

\paragraph{TTS Orchestrator}
Coordina la generazione audio tramite modelli di sintesi vocale.  
Responsabilità:
\begin{itemize}
    \item Conversione dei chunk narrativi in audio tramite provider esterni.
    \item Gestione di parametri vocali: timbro, tono, velocità, stile.
    \item Rendering sia in modalità ``Fast'' che ``High Quality''.
    \item Ottimizzazione dei costi mediante caching e pre-rendering.
    \item Versioning delle voci utilizzate per ogni episodio.
\end{itemize}

\paragraph{Streaming Orchestrator / Player Backend}
Il modulo di orchestrazione della riproduzione audio.  
Responsabilità:
\begin{itemize}
    \item Creazione della timeline audio che combina musica esterna e parlato AI.
    \item Sincronizzazione tra due flussi indipendenti:
    \begin{itemize}
        \item musica da Spotify/Apple Music (riprodotta localmente dal dispositivo dell’utente);
        \item voce AI proveniente dal CDN di Radio-AI.
    \end{itemize}
    \item Gestione dei segnali di controllo (play, pausa, skip, aggiornamento prompt).
    \item Supporto alla riproduzione sincronizzata nelle Radio Rooms via WebSocket.
    \item Monitoraggio del buffering, del drift temporale e delle latenze.
\end{itemize}

---

\subsection{Radio Rooms Service}

Le Radio Rooms consentono l'ascolto collettivo di una puntata Radio-AI in tempo reale, garantendo che tutti i partecipanti, anche da dispositivi e posizioni diverse, ascoltino esattamente lo stesso contenuto nello stesso momento. Questo modulo rappresenta uno dei componenti più distintivi della piattaforma, poiché introduce un livello di interazione sociale che combina tecnologia di sincronizzazione a bassa latenza e gestione distribuita degli utenti.

\paragraph{Responsabilità principali}
\begin{itemize}
    \item \textbf{Creazione delle stanze}: ogni episodio generato può opzionalmente creare una Radio Room associata.
    \item \textbf{Join/Leave}: gestione dell’ingresso e dell’uscita degli utenti, incluse riconnessioni e handover di stato.
    \item \textbf{Sincronizzazione temporale}: mantenimento di un timestamp di riferimento condiviso per tutti i client con tolleranza di pochi millisecondi.
    \item \textbf{Host controls}: l’utente host può controllare avvio, pausa, skip, modifica del prompt e impostazioni della stanza.
    \item \textbf{Prompt update broadcasting}: se l'host modifica il prompt, i nuovi chunk vengono rigenerati e sincronizzati su tutti i partecipanti.
    \item \textbf{Messaggi WebSocket}: trasmissione di heartbeat, timestamp update, richieste di sincronizzazione e notifiche lato server.
\end{itemize}

\paragraph{Architettura}
Il servizio si basa su:
\begin{itemize}
    \item un \textbf{WebSocket Hub} che gestisce connessioni persistenti;
    \item un \textbf{Sync Server} che aggiorna e distribuisce il clock della stanza;
    \item un \textbf{Session Store} (tipicamente Redis) per mantenere stato, ruoli e offset temporali;
    \item un \textbf{Drift Controller} per correggere differenze di latenza tra dispositivi.
\end{itemize}

\begin{figure}[H]
    \centering
    \includegraphics[width=0.8\textwidth]{fig/radio-rooms-sync.png}
    \caption{Radio Rooms: architettura di sincronizzazione multi–device e gestione real-time.}
    \label{fig:radio-rooms}
\end{figure}

\paragraph{Benefici architetturali}
\begin{itemize}
    \item Scalabilità lineare grazie alla separazione tra WebSocket Hub e servizi di generazione contenuti.
    \item Robustezza in caso di lag o variazioni di rete tramite resync progressivo.
    \item Compatibilità con musica riprodotta localmente tramite Spotify/Apple Music e parlato via CDN.
\end{itemize}

% -------------------------------------------------------------

\subsection{Social Layer \& Content Management}

Il Social Layer trasforma Radio-AI da semplice generatore di audio personalizzato a una piattaforma completa di contenuti condivisibili, esplorabili e riutilizzabili.  
Questo modulo gestisce identità utente, librerie di episodi, Quote Cards e meccanismi di condivisione e discovery.

\paragraph{Componenti principali}

\begin{itemize}
    \item \textbf{User Profile Service}  
    Gestisce profili pubblici o privati, incluse preferenze, avatar, impostazioni vocali e storici di utilizzo.

    \item \textbf{Episode Library Service}  
    Archivia tutte le puntate generate, incluse:
    \begin{itemize}
        \item prompt originale e prompt rigenerato dall’AI, utilizzato come “descrizione ufficiale” dell’episodio per catalogazione, ricerca e condivisione;
        \item trascrizione completa;
        \item tag tematici generati automaticamente;
        \item metadati (durata, mood, stile, rubriche incluse).
    \end{itemize}

    Il processo di rigenerazione del prompt trasforma input anche lunghi o colloquiali in una descrizione chiara, leggibile e in stile “episodio podcast”, facilitando l’anteprima della puntata e l’indicizzazione semantica all’interno del motore di ricerca.

    \item \textbf{Quote Cards Generator}  
    Trasforma selezioni di testo dalla trascrizione in card grafiche condivisibili, con possibilità di personalizzare layout, font e colori.

    \item \textbf{Public Episode Feed}  
    Consente agli utenti di scoprire puntate pubbliche generate da altri, ordinabili per:
    \begin{itemize}
        \item trend;
        \item tema;
        \item mood;
        \item artista principale;
        \item creator.
    \end{itemize}

    \item \textbf{Search \& Recommendation Engine}  
    Motore di ricerca su:
    \begin{itemize}
        \item titolo dell’episodio;
        \item prompt rigenerato;
        \item tag tematici;
        \item contenuti parlati.
    \end{itemize}
    Utilizza embedding semantici per facilitare la discovery dei contenuti.
\end{itemize}

\paragraph{Funzionalità secondarie}
\begin{itemize}
    \item Like, salvataggi e condivisione puntate.
    \item Esportazione card e link univoci per ascolto.
    \item Cronologia e Re–generation: ogni utente può recuperare le proprie puntate passate dalla cronologia personale e rigenerarle non in modo identico, ma mantenendo tema e formato di base, ampliando o aggiornando i contenuti; inoltre può rigenerare puntate pubbliche altrui adattandole al proprio stile.
\end{itemize}

\paragraph{Obiettivi del Social Layer}
\begin{itemize}
    \item aumentare la retention attraverso la costruzione di un ecosistema di contenuti;
    \item favorire comportamenti di condivisione e partecipazione collettiva;
    \item creare valore anche oltre l’ascolto (es. Quote Cards come asset social).
\end{itemize}

% -------------------------------------------------------------

\subsection{Radio Ads Manager}

Il \textbf{Radio Ads Manager} è l’infrastruttura dedicata alle aziende che desiderano inserire contenuti pubblicitari all’interno delle dirette generate.  
A differenza dei sistemi tradizionali basati su targeting demografico, Radio-AI introduce un nuovo concetto di \textbf{targeting tematico ed emozionale}.

\paragraph{Responsabilità principali}
\begin{itemize}
    \item Creazione e configurazione delle campagne (obiettivo, budget, durata).
    \item Targeting basato su temi, atmosfere, mood e generi musicali.
    \item Generazione automatica dello spot:
    \begin{itemize}
        \item script scritto dall’AI;
        \item voce sintetica in linea con la campagna;
        \item eventuale integrazione narrativa con la puntata.
    \end{itemize}
    \item Posizionamento dinamico dello spot nelle puntate coerenti.
    \item Reportistica dettagliata (impressions, completion rate, contesti tematici raggiunti).
\end{itemize}

\paragraph{Modalità di inserimento}
Gli spot possono essere inseriti in:
\begin{itemize}
    \item intro o outro della puntata;
    \item pre–brano o post–brano;
    \item rubriche sponsorizzate;
    \item transizioni narrative.
\end{itemize}

\paragraph{Architettura Funzionale}
\begin{itemize}
    \item \textbf{Campaign Manager}: archivio e gestione delle campagne attive.
    \item \textbf{Target Matching Engine}: seleziona le puntate coerenti con il tema.
    \item \textbf{Insertion Engine}: decide il placement ottimale all’interno dell’episodio.
    \item \textbf{Ad Rendering Service}: genera spot audio via TTS.
    \item \textbf{Insights \& Analytics Service}: calcola metriche e produce report.
\end{itemize}

\begin{figure}[H]
    \centering
    \includegraphics[width=0.8\textwidth]{fig/radio-ads-manager.png}
    \caption{Targeting tematico e inserimento contestuale degli spot all’interno delle puntate.}
    \label{fig:ads-manager}
\end{figure}
% ---------------------------------
\section{Data Model}

\subsection{Core Entities}

Il modello dati di Radio-AI è centrato su alcune entità principali che rappresentano utenti, contenuti generati, strutture di ascolto condiviso e campagne pubblicitarie.  
In questa sezione descriviamo le entità chiave a livello concettuale.

\paragraph{User}
Rappresenta un utilizzatore della piattaforma (ascoltatore, creator, host di Radio Rooms, advertiser in caso B2B).
\begin{itemize}
    \item Identità (ID utente, provider OAuth, email opzionale).
    \item Preferenze (lingua, stile vocale preferito, temi ricorrenti).
    \item Stato account (free, premium, advertiser).
    \item Metadati social (bio, visibilità del profilo, followers/following).
\end{itemize}

\paragraph{Prompt}
Rappresenta la richiesta testuale che dà origine a una puntata.
\begin{itemize}
    \item Prompt originale inserito dall’utente.
    \item Prompt rigenerato (versione pulita in stile “descrizione episodio”).
    \item Intenti estratti (tema, rubriche, mood, preferenze musicali).
\end{itemize}

\paragraph{Episode}
Rappresenta una singola diretta generata da Radio-AI.
\begin{itemize}
    \item Riferimento all’utente creatore.
    \item Riferimento al prompt associato.
    \item Durata e stato (draft, generato, pubblicato, archiviato).
    \item Metadati (tema, mood, stile dello speaker).
    \item Collegamento a playlist e chunk audio.
\end{itemize}

\paragraph{Playlist}
Rappresenta la selezione musicale utilizzata in una puntata.
\begin{itemize}
    \item Origine (utente, AI, servizio esterno).
    \item Lista ordinata di brani (ID esterni tipo Spotify/Apple).
    \item Vincoli o annotazioni (brani obbligatori, brani esclusi).
\end{itemize}

\paragraph{Chunk}
Rappresenta un blocco di parlato o contenuto audio generato dall’AI.
\begin{itemize}
    \item Tipo di chunk (intro, collegamento, rubrica, spot, outro).
    \item Testo sorgente (script).
    \item File audio TTS generato (location su storage/CDN).
    \item Posizione temporale relativa all’episodio.
    \item Versioni (in caso di rigenerazioni successive).
\end{itemize}

\paragraph{Room}
Rappresenta una Radio Room attiva o passata.
\begin{itemize}
    \item Riferimento all’episodio di base.
    \item Host e lista dei partecipanti.
    \item Stato della stanza (attiva, chiusa).
    \item Timestamp di riferimento e offset per sincronizzazione.
\end{itemize}

\paragraph{QuoteCard}
Rappresenta una card grafica generata a partire da una frase o estratto della puntata.
\begin{itemize}
    \item Riferimento all’episodio e alla posizione nel testo.
    \item Testo della citazione.
    \item Template grafico utilizzato.
    \item URL dell’immagine generata.
\end{itemize}

\paragraph{Tag}
Rappresenta un’etichetta tematica associata ad episodi, prompt o utenti.
\begin{itemize}
    \item Tipo di tag (tema, mood, genere musicale, contesto).
    \item Valore normalizzato (es. \texttt{viaggi}, \texttt{nostalgia-anni-90}).
\end{itemize}

\paragraph{Campaign (Ads)}
Rappresenta una campagna pubblicitaria configurata nel Radio Ads Manager.
\begin{itemize}
    \item Advertiser (user o account aziendale).
    \item Obiettivo (vendite, traffico, awareness).
    \item Budget totale e finestra temporale.
    \item Parametri di targeting tematico.
\end{itemize}

\paragraph{AdSpot}
Rappresenta un singolo inserimento di uno spot all’interno di una puntata.
\begin{itemize}
    \item Riferimento alla campagna.
    \item Episodio e posizione (pre–brano, post–brano, rubrica).
    \item Chunk audio dello spot.
    \item Metriche (impressions, completamento, ascolto medio).
\end{itemize}

% -------------------------------------------------------------

\subsection{Example Schema}

Questa sezione fornisce una vista semplificata del modello dati in forma tabellare, utile come base per un ER diagram più dettagliato.

\subsubsection*{Tabella \texttt{users}}

\begin{table}[H]
    \centering
    \begin{tabular}{llp{7cm}}
        \toprule
        \textbf{Campo} & \textbf{Tipo} & \textbf{Descrizione} \\
        \midrule
        id & UUID & Identificatore univoco utente \\
        auth\_provider & VARCHAR & Provider OAuth (spotify, apple, google, ecc.) \\
        auth\_subject & VARCHAR & ID esterno presso il provider \\
        email & VARCHAR (NULL) & Email utente (opzionale) \\
        display\_name & VARCHAR & Nome visualizzato nel profilo \\
        account\_type & ENUM & \{free, premium, advertiser\} \\
        created\_at & TIMESTAMP & Data creazione account \\
        updated\_at & TIMESTAMP & Ultimo aggiornamento \\
        \bottomrule
    \end{tabular}
    \caption{Schema semplificato della tabella \texttt{users}.}
\end{table}

\subsubsection*{Tabella \texttt{prompts}}

\begin{table}[H]
    \centering
    \begin{tabular}{llp{7cm}}
        \toprule
        \textbf{Campo} & \textbf{Tipo} & \textbf{Descrizione} \\
        \midrule
        id & UUID & Identificatore univoco prompt \\
        user\_id & UUID & Riferimento a \texttt{users.id} \\
        raw\_text & TEXT & Prompt originale inserito dall’utente \\
        normalized\_text & TEXT & Prompt rigenerato/normalizzato \\
        intents\_json & JSONB & Intent estratti (tema, mood, rubriche, musica) \\
        created\_at & TIMESTAMP & Data creazione \\
        \bottomrule
    \end{tabular}
    \caption{Schema semplificato della tabella \texttt{prompts}.}
\end{table}

\subsubsection*{Tabella \texttt{episodes}}

\begin{table}[H]
    \centering
    \begin{tabular}{llp{7cm}}
        \toprule
        \textbf{Campo} & \textbf{Tipo} & \textbf{Descrizione} \\
        \midrule
        id & UUID & Identificatore univoco episodio \\
        user\_id & UUID & Creatore dell’episodio (\texttt{users.id}) \\
        prompt\_id & UUID & Prompt di origine (\texttt{prompts.id}) \\
        title & VARCHAR & Titolo generato dall’AI \\
        description & TEXT & Descrizione estesa (prompt rigenerato) \\
        duration\_sec & INT & Durata prevista dell’episodio in secondi \\
        status & ENUM & \{draft, generated, published, archived\} \\
        metadata\_json & JSONB & Tema, mood, stile, rubriche attive \\
        created\_at & TIMESTAMP & Data generazione \\
        \bottomrule
    \end{tabular}
    \caption{Schema semplificato della tabella \texttt{episodes}.}
\end{table}

\subsubsection*{Tabella \texttt{chunks}}

\begin{table}[H]
    \centering
    \begin{tabular}{llp{7cm}}
        \toprule
        \textbf{Campo} & \textbf{Tipo} & \textbf{Descrizione} \\
        \midrule
        id & UUID & Identificatore univoco chunk \\
        episode\_id & UUID & Episodio di appartenenza (\texttt{episodes.id}) \\
        chunk\_index & INT & Ordine del chunk nella timeline \\
        chunk\_type & ENUM & \{intro, bridge, rubric, ad, outro, other\} \\
        script\_text & TEXT & Testo sorgente del parlato \\
        audio\_url & VARCHAR & URL al file audio TTS nel CDN \\
        start\_offset\_sec & INT & Offset (secondi) rispetto all’inizio dell’episodio \\
        version & INT & Numero di versione del chunk (per rigenerazioni) \\
        \bottomrule
    \end{tabular}
    \caption{Schema semplificato della tabella \texttt{chunks}.}
\end{table}

\subsubsection*{Tabella \texttt{rooms}}

\begin{table}[H]
    \centering
    \begin{tabular}{llp{7cm}}
        \toprule
        \textbf{Campo} & \textbf{Tipo} & \textbf{Descrizione} \\
        \midrule
        id & UUID & Identificatore univoco stanza \\
        episode\_id & UUID & Episodio di base (\texttt{episodes.id}) \\
        host\_user\_id & UUID & Utente host (\texttt{users.id}) \\
        status & ENUM & \{active, closed\} \\
        current\_timestamp\_sec & INT & Timestamp condiviso dell’episodio \\
        created\_at & TIMESTAMP & Creazione stanza \\
        closed\_at & TIMESTAMP (NULL) & Chiusura stanza, se presente \\
        \bottomrule
    \end{tabular}
    \caption{Schema semplificato della tabella \texttt{rooms}.}
\end{table}

\subsubsection*{Tabelle \texttt{campaigns} e \texttt{ad\_spots}}

\begin{table}[H]
    \centering
    \begin{tabular}{llp{7cm}}
        \toprule
        \textbf{Campo} & \textbf{Tipo} & \textbf{Descrizione} \\
        \midrule
        id & UUID & Identificatore univoco campagna \\
        advertiser\_id & UUID & Utente/azienda (\texttt{users.id}) \\
        objective & ENUM & \{sales, traffic, awareness\} \\
        budget\_total & NUMERIC & Budget totale \\
        start\_date & DATE & Inizio campagna \\
        end\_date & DATE & Fine campagna \\
        targeting\_json & JSONB & Temi, mood, generi, fasce orarie target \\
        status & ENUM & \{draft, active, paused, completed\} \\
        \bottomrule
    \end{tabular}
    \caption{Schema semplificato della tabella \texttt{campaigns}.}
\end{table}

\begin{table}[H]
    \centering
    \begin{tabular}{llp{7cm}}
        \toprule
        \textbf{Campo} & \textbf{Tipo} & \textbf{Descrizione} \\
        \midrule
        id & UUID & Identificatore univoco spot \\
        campaign\_id & UUID & Campagna associata (\texttt{campaigns.id}) \\
        episode\_id & UUID & Episodio in cui lo spot è inserito \\
        chunk\_id & UUID & Chunk audio dello spot (\texttt{chunks.id}) \\
        placement\_type & ENUM & \{intro, midroll, preroll\_song, postroll\_song, outro\} \\
        impressions & INT & Numero di ascolti avvenuti \\
        completion\_rate & NUMERIC & Percentuale di ascolti completati \\
        \bottomrule
    \end{tabular}
    \caption{Schema semplificato della tabella \texttt{ad\_spots}.}
\end{table}

Questi schemi non rappresentano un modello definitivo, ma costituiscono una base concreta da cui partire per la progettazione di un database relazionale (ad esempio PostgreSQL) o per l’estrazione di un ER diagram formale.
% ---------------------------------
\section{AI Components}

\subsection{Prompt Interpretation \& Intent Extraction}

La pipeline di interpretazione del prompt è uno degli elementi centrali di Radio-AI.  
L’obiettivo è trasformare un input testuale libero (spesso colloquiale, ambiguo o incompleto) in un set strutturato di istruzioni operative da cui generare un episodio coerente.

\paragraph{Prompt Interpretation}

Questa fase utilizza un Large Language Model (LLM) specializzato nella comprensione semantica e nella classificazione dei contenuti.  
Responsabilità principali:
\begin{itemize}
    \item Analisi del linguaggio naturale per estrarre:
    \begin{itemize}
        \item tema dell’episodio;
        \item tono e stile desiderato dello speaker;
        \item rubriche o formati da includere;
        \item richieste culturali (curiosità, aneddoti, approfondimenti);
        \item preferenze musicali esplicite (artisti, epoche, generi);
        \item frequenza o distribuzione della musica nella puntata.
    \end{itemize}

    \item Identificazione di richieste implicite o contestuali:
    \begin{itemize}
        \item mood dell’utente (es. “sono agitato prima di un appuntamento”);
        \item obiettivi (“dare motivazione”, “riposare”, “studiare”);
        \item intenzioni stilistiche non dichiarate ma deducibili (es. ritmo lento vs dinamico).
    \end{itemize}

    \item Normalizzazione del prompt in una struttura JSON interpretabile dal sistema.
\end{itemize}

\paragraph{Intent Extraction}

Dopo l’interpretazione preliminare, un modulo dedicato di \textbf{Intent Extractor} verifica, pulisce e consolida gli intenti.  
Responsabilità:
\begin{itemize}
    \item disambiguazione dei parametri conflittuali (es. “solo 1 canzone” vs “canzoni frequenti”);
    \item completamento automatico di parametri mancanti usando euristiche;
    \item validazione delle richieste rispetto ai limiti della piattaforma;
    \item normalizzazione delle rubriche in categorie standardizzate.
\end{itemize}

Il risultato finale è un set di istruzioni strutturate che guidano tutte le fasi successive della pipeline.

% -------------------------------------------------------------

\subsection{Script Generation}

Lo \textbf{Script Generator} è responsabile della creazione della narrazione completa della puntata.  
La generazione avviene in modo incrementale e strutturato, per garantire coerenza, modularità e possibilità di rigenerazione parziale.

\paragraph{Struttura dello script}

Uno script è tipicamente suddiviso in:
\begin{itemize}
    \item \textbf{Intro}: apertura narrativa, definizione del mood.
    \item \textbf{Bridges}: collegamenti tra i brani.
    \item \textbf{Rubriche}: approfondimenti culturali, storici, musicali o tematici.
    \item \textbf{Elementi di intrattenimento}: telefonate finte, dialoghi immaginari, sketch.
    \item \textbf{Segmenti sponsorizzati}: se una campagna pubblicitaria è attiva.
    \item \textbf{Outro}: chiusura coerente con il tema.
\end{itemize}

\paragraph{Chunked Generation}

Per ottimizzare la latenza e permettere la modifica live del prompt, la generazione dello script segue il paradigma:
\[
\text{episodio} = \sum_{i=1}^{N} \text{chunk}_i
\]

Ogni chunk contiene:
\begin{itemize}
    \item testo del parlato (30–90s);
    \item tipo del segmento (intro, collegamento, rubrica, fiction, ecc.);
    \item riferimenti alla musica circostante;
    \item eventuali entità culturali verificate.
\end{itemize}

\paragraph{Contenuti reali vs contenuti fittizi}

Lo Script Generator applica regole rigide per rispettare la distinzione di Radio-AI tra ciò che è reale e ciò che è inventato:

\begin{itemize}
    \item \textbf{curiosità, aneddoti e contenuti culturali}: devono provenire dal Fact-Checking Engine e vengono inclusi solo se veri, verificabili e coerenti con le fonti;
    \item \textbf{contenuti potenzialmente non verificabili}: vengono omessi oppure riscritti in forma certa, evitando qualsiasi ambiguità;
    \item \textbf{contenuti di pura finzione}: (telefonate immaginarie, sketch, momenti narrativi) sono ammessi solo come elementi di intrattenimento radiofonico e non vengono mai presentati come fatti, news o informazioni culturali.
\end{itemize}

% -------------------------------------------------------------

\subsection{Fact-Checking Engine}

Il Fact-Checking Engine garantisce che tutte le curiosità, le informazioni storiche, le citazioni culturali e gli aneddoti presenti nella puntata siano:
\begin{itemize}
    \item veri,
    \item verificabili,
    \item tracciabili a fonti affidabili.
\end{itemize}

\paragraph{Pipeline RAG (Retrieval-Augmented Generation)}

La pipeline utilizza un sistema RAG con i seguenti step:

\begin{enumerate}
    \item \textbf{Extraction}: dallo script vengono estratte le frasi che richiedono verifica (nomi propri, date, eventi, dischi, premi, luoghi).
    \item \textbf{Retrieval}: interrogazione di fonti affidabili quali:
    \begin{itemize}
        \item database musicali (MusicBrainz, Discogs);
        \item enciclopedie strutturate (Wikipedia + Wikidata);
        \item dataset culturali e storici certificati;
        \item knowledge graph proprietari.
    \end{itemize}
    \item \textbf{Alignment}: confronto tra il contenuto generato e i dati recuperati.
    \item \textbf{Rewrite/Omission}: contenuti non verificabili vengono:
    \begin{itemize}
        \item riscritti in forma corretta,
        \item semplificati,
        \item oppure eliminati del tutto.
    \end{itemize}
\end{enumerate}

\paragraph{Caching}

Il Fact-Checking Engine implementa:
\begin{itemize}
    \item caching di artisti, album, eventi e date frequentemente richiesti;
    \item caching di rubriche popolari (es. “curiosità su David Bowie”).
\end{itemize}

Questo riduce latenza e costi operativi.

% -------------------------------------------------------------

\subsection{Text-to-Speech \& Voice Personalization}

Il modulo TTS è responsabile della conversione dello script testuale in voce naturale, mantenendo coerenza stilistica, intonazione e ritmo.

\paragraph{Parametri vocali}

Ogni voce può essere configurata attraverso:
\begin{itemize}
    \item \textbf{timbro}: maschile, femminile, neutro;
    \item \textbf{tono}: caldo, profondo, squillante, intimo, ironico;
    \item \textbf{ritmo}: lento, medio, veloce;
    \item \textbf{stile}: radiofonico, documentaristico, podcast, DJ, storyteller;
    \item \textbf{postura vocale}: drammatica, empatica, energica, calma.
\end{itemize}

\paragraph{Provider TTS}

Radio-AI può integrare diversi provider:
\begin{itemize}
    \item Google/DeepMind TTS;
    \item Microsoft Azure Neural VOX;
    \item ElevenLabs;
    \item soluzioni self-hosted per voci proprietarie.
\end{itemize}

\paragraph{Voci Premium e Licenze}

Il sistema supporta:
\begin{itemize}
    \item voci “ispirate” a personaggi famosi (senza imitazione perfetta);
    \item voci firmate tramite accordi commerciali;
    \item voci caricabili dall’utente (user-generated voice cloning), con necessità di consenso esplicito.
\end{itemize}

\paragraph{Rendering}

Il TTS Orchestrator supporta:
\begin{itemize}
    \item \textbf{Fast Mode}: bassa latenza per avvio immediato della puntata;
    \item \textbf{HQ Mode}: rendering di qualità più alta per download o salvataggio definitivo;
    \item \textbf{Chunked Rendering}: generazione incrementale per supporto alla modifica in tempo reale del prompt.
\end{itemize}
% ---------------------------------
\section{Feature Flow Diagrams}

\subsection{Flow: Generate Episode with User Playlist (Mode A)}

In Modalità A l'utente fornisce una playlist esistente (es. da Spotify o Apple Music) e Radio-AI costruisce la diretta utilizzando esclusivamente quella selezione musicale.

\paragraph{Passi principali}

\begin{enumerate}
    \item \textbf{Selezione playlist}
    \begin{enumerate}
        \item L'utente effettua il login e collega il proprio account Spotify/Apple.
        \item Dall'interfaccia, seleziona una playlist esistente.
        \item Il client invia all'API Gateway l'ID della playlist esterna e il prompt opzionale.
    \end{enumerate}

    \item \textbf{Interpretazione del prompt}
    \begin{enumerate}
        \item L'API Gateway inoltra il prompt al \textit{Prompt Interpreter Service}.
        \item Il servizio estrae tema, mood, rubriche, stile dello speaker e parametri di durata.
        \item Gli intenti vengono validati dall'\textit{Intent Extractor} e serializzati in JSON.
    \end{enumerate}

    \item \textbf{Integrazione playlist utente}
    \begin{enumerate}
        \item Il \textit{Playlist Service} interroga l'API di Spotify/Apple per recuperare la lista di brani.
        \item I metadati musicali (titolo, artista, durata, ID) vengono normalizzati e collegati all'episodio.
        \item La sequenza dei brani è considerata fissa (salvo modalità shuffle opzionale richiesta dall'utente).
    \end{enumerate}

    \item \textbf{Generazione script}
    \begin{enumerate}
        \item Il \textit{Script Generator Service} riceve:
        \begin{itemize}
            \item JSON degli intenti;
            \item metadati della playlist;
            \item eventuali vincoli (es. parlare poco, focus su curiosità, ecc.).
        \end{itemize}
        \item Viene costruita una scaletta con:
        \begin{itemize}
            \item intro generale;
            \item blocchi di parlato tra i brani;
            \item rubriche e curiosità culturali, verificate tramite \textit{Fact-Checking Engine};
            \item outro conclusiva.
        \end{itemize}
        \item Lo script viene suddiviso in chunk.
    \end{enumerate}

    \item \textbf{Rendering audio}
    \begin{enumerate}
        \item Ogni chunk viene inviato al \textit{TTS Orchestrator}.
        \item Viene generato l'audio in modalità ``Fast'' per l'avvio rapido della puntata.
        \item In parallelo può essere avviato il rendering ``HQ''.
    \end{enumerate}

    \item \textbf{Playback}
    \begin{enumerate}
        \item Il \textit{Streaming Orchestrator} genera la timeline dell'episodio.
        \item Il client:
        \begin{itemize}
            \item riproduce la musica via Spotify/Apple;
            \item riproduce i chunk audio via CDN.
        \end{itemize}
        \item L'utente ascolta la diretta radiofonica basata interamente sulla propria playlist.
    \end{enumerate}
\end{enumerate}

% -------------------------------------------------------------

\subsection{Flow: Generate Episode without Playlist (Mode B)}

In Modalità B l'utente non seleziona alcuna playlist: è Radio-AI a generare la selezione musicale, lasciando comunque la possibilità di inserire brani specifici prima dell'avvio definitivo.

\paragraph{Passi principali}

\begin{enumerate}
    \item \textbf{Input del prompt}
    \begin{enumerate}
        \item L'utente inserisce un prompt (anche molto generico, es. ``puntata motivazionale per allenamento'').
        \item Il client invia il prompt all'API Gateway.
    \end{enumerate}

    \item \textbf{Interpretazione e intenti}
    \begin{enumerate}
        \item Il \textit{Prompt Interpreter Service} estrae:
        \begin{itemize}
            \item tema e mood;
            \item stile vocale;
            \item preferenze musicali implicite (es. generi energici, epoche, artisti menzionati).
        \end{itemize}
        \item L'\textit{Intent Extractor} produce un JSON normalizzato di intenti.
    \end{enumerate}

    \item \textbf{Generazione playlist AI-driven}
    \begin{enumerate}
        \item Il \textit{Playlist Service} utilizza il JSON degli intenti per selezionare brani coerenti:
        \begin{itemize}
            \item tramite raccomandazioni basate su generi, mood e BPM;
            \item tramite integrazioni con API esterne (es. Spotify recommendations).
        \end{itemize}
        \item Viene generata una playlist proposta per la puntata.
    \end{enumerate}

    \item \textbf{Richiesta di brani aggiuntivi}
    \begin{enumerate}
        \item Prima di procedere, il sistema chiede all'utente:
        \begin{quote}
            ``Vuoi aggiungere qualche brano specifico alla puntata?''
        \end{quote}
        \item Se l'utente aggiunge brani:
        \begin{itemize}
            \item questi vengono integrati nella playlist AI;
            \item eventuali duplicati o conflitti vengono risolti dal \textit{Playlist Service}.
        \end{itemize}
        \item Se l'utente non aggiunge nulla:
        \begin{itemize}
            \item la playlist AI viene utilizzata così com'è.
        \end{itemize}
    \end{enumerate}

    \item \textbf{Generazione script e rendering}
    \begin{enumerate}
        \item Il flusso prosegue come in Modalità A:
        \begin{itemize}
            \item generazione script a chunk;
            \item fact-checking per contenuti culturali;
            \item rendering TTS;
            \item creazione timeline di playback.
        \end{itemize}
    \end{enumerate}
\end{enumerate}

% -------------------------------------------------------------

\subsection{Flow: Live Prompt Update During Episode}

Questa funzionalità permette all'utente di modificare il prompt mentre l'episodio è in corso, senza dover ricominciare da capo.

\paragraph{Passi principali}

\begin{enumerate}
    \item \textbf{Pause \& Snapshot}
    \begin{enumerate}
        \item L'utente mette in pausa la puntata.
        \item Il \textit{Session Service} registra:
        \begin{itemize}
            \item timestamp corrente;
            \item chunk in riproduzione;
            \item contesto narrativo attuale.
        \end{itemize}
    \end{enumerate}

    \item \textbf{Aggiornamento prompt}
    \begin{enumerate}
        \item L'utente modifica il prompt (es. cambia tema, mood, stile dello speaker).
        \item Il nuovo prompt viene inviato al \textit{Prompt Interpreter Service}.
        \item Viene generato un set aggiornato di intenti.
    \end{enumerate}

    \item \textbf{Rigenerazione dei chunk futuri}
    \begin{enumerate}
        \item Il \textit{Script Generator Service} mantiene invariati:
        \begin{itemize}
            \item i chunk già riprodotti;
            \item i segmenti già confermati dall'utente.
        \end{itemize}
        \item Rigenera solo i chunk successivi al timestamp corrente, allineandoli ai nuovi intenti.
        \item I nuovi chunk vengono passati al \textit{TTS Orchestrator} per il rendering.
    \end{enumerate}

    \item \textbf{Resume}
    \begin{enumerate}
        \item Alla pressione di ``Play'', il \textit{Streaming Orchestrator} riprende dal punto in cui la puntata era stata interrotta.
        \item I nuovi chunk sostituiscono quelli precedentemente pianificati per la parte restante dell'episodio.
    \end{enumerate}
\end{enumerate}

% -------------------------------------------------------------

\subsection{Flow: Radio Room Creation \& Join}

Questo flusso descrive la creazione di una Radio Room e la partecipazione di più utenti all'ascolto sincronizzato.

\paragraph{Passi principali}

\begin{enumerate}
    \item \textbf{Creazione stanza}
    \begin{enumerate}
        \item L'utente seleziona un episodio e sceglie ``Crea Radio Room''.
        \item Il client invia una richiesta al \textit{Radio Rooms Service}.
        \item Viene creata una nuova entry \texttt{room} con:
        \begin{itemize}
            \item host\_user\_id;
            \item episode\_id;
            \item timestamp iniziale (0 o punto specifico).
        \end{itemize}
        \item Viene generato un link/un codice di invito.
    \end{enumerate}

    \item \textbf{Join dei partecipanti}
    \begin{enumerate}
        \item Gli altri utenti cliccano sul link o inseriscono il codice.
        \item Ogni client apre una connessione WebSocket verso il WebSocket Hub.
        \item Il \textit{Radio Rooms Service} registra i partecipanti nella stanza.
    \end{enumerate}

    \item \textbf{Sincronizzazione}
    \begin{enumerate}
        \item Il \textit{Sync Server} mantiene un timestamp condiviso per la stanza.
        \item I client ricevono periodicamente:
        \begin{itemize}
            \item timestamp di riferimento;
            \item eventuali correzioni (drift correction).
        \end{itemize}
        \item Ogni client allinea:
        \begin{itemize}
            \item playback della musica tramite Spotify/Apple;
            \item playback dei chunk vocali via CDN.
        \end{itemize}
    \end{enumerate}

    \item \textbf{Controlli dell’host}
    \begin{enumerate}
        \item L'host può:
        \begin{itemize}
            \item mettere in pausa;
            \item skippare;
            \item modificare il prompt.
        \end{itemize}
        \item Queste azioni vengono propagate via WebSocket a tutti i partecipanti.
    \end{enumerate}
\end{enumerate}

% -------------------------------------------------------------

\subsection{Flow: Ad Campaign Creation \& Placement}

Questo flusso descrive il ciclo di vita di una campagna pubblicitaria nel Radio Ads Manager, dal setup all’inserimento degli spot nelle puntate.

\paragraph{Passi principali}

\begin{enumerate}
    \item \textbf{Creazione campagna}
    \begin{enumerate}
        \item Un advertiser accede all'interfaccia del Radio Ads Manager.
        \item Seleziona un obiettivo (vendite, traffico, awareness).
        \item Definisce budget totale, durata e limiti giornalieri opzionali.
    \end{enumerate}

    \item \textbf{Targeting tematico}
    \begin{enumerate}
        \item L’advertiser seleziona:
        \begin{itemize}
            \item temi (es. viaggi, sport, relax);
            \item mood (energetico, notturno, nostalgico);
            \item generi musicali predominanti;
            \item fasce orarie preferite.
        \end{itemize}
        \item Queste informazioni vengono salvate in \texttt{targeting\_json}.
    \end{enumerate}

    \item \textbf{Generazione contenuto}
    \begin{enumerate}
        \item L’advertiser può:
        \begin{itemize}
            \item caricare uno script esistente;
            \item oppure chiedere all'AI di generare lo script a partire da un prompt.
        \end{itemize}
        \item Lo script viene inviato al \textit{TTS Orchestrator} per la creazione dello spot audio.
    \end{enumerate}

    \item \textbf{Attivazione e matching}
    \begin{enumerate}
        \item La campagna passa in stato \texttt{active}.
        \item Il \textit{Target Matching Engine} monitora gli episodi in generazione:
        \begin{itemize}
            \item confronta tema, tag e mood dell'episodio con il targeting della campagna;
            \item seleziona solo le puntate coerenti.
        \end{itemize}
    \end{enumerate}

    \item \textbf{Placement nello show}
    \begin{enumerate}
        \item L'\textit{Insertion Engine} decide i punti di inserimento:
        \begin{itemize}
            \item intro, mid-roll, pre/post brano, outro;
            \item rubriche sponsorizzate.
        \end{itemize}
        \item Vengono creati \texttt{ad\_spots} associati ai relativi episodi e chunk.
    \end{enumerate}

    \item \textbf{Reporting}
    \begin{enumerate}
        \item Durante l’ascolto, il sistema registra:
        \begin{itemize}
            \item impressions;
            \item completion rate;
            \item contesti tematici in cui lo spot è stato riprodotto.
        \end{itemize}
        \item L'\textit{Insights \& Analytics Service} espone report aggregati all'advertiser.
    \end{enumerate}
\end{enumerate}
% ---------------------------------
\section{Technology Stack}

Il seguente stack tecnologico rappresenta una proposta realistica, scalabile e moderna per implementare Radio-AI in produzione.  
L’architettura è modulare e supporta l’evoluzione progressiva del prodotto (nuove feature, multi-provider AI, crescita del traffico).

% -------------------------------------------------------------

\subsection{Backend}

Il backend di Radio-AI beneficia di un approccio a microservizi, con componenti indipendenti e coordinati tramite API interne.

\paragraph{Linguaggi e Framework}
\begin{itemize}
    \item \textbf{Node.js (TypeScript)} — ideale per API ad alta concorrenza, WebSocket, orchestrazione.
    \item \textbf{Python} — utilizzato per i moduli AI intensivi (Prompt Interpreter, RAG, Fact-Checking).
    \item \textbf{FastAPI / Flask} — per i servizi di AI + Fact-Checking.
    \item \textbf{gRPC o REST JSON} — per comunicazioni tra microservizi.
\end{itemize}

\paragraph{Architettura}
\begin{itemize}
    \item \textbf{Microservizi indipendenti}: API Gateway, Session Service, Playlist Service, Script Generator, TTS Orchestrator, Streaming Orchestrator, Radio Rooms Service, Ads Manager Backend.
    \item \textbf{Event-driven patterns} tramite code (es. RabbitMQ / Kafka) per:
    \begin{itemize}
        \item generazione chunk,
        \item rendering TTS,
        \item inserzioni,
        \item analisi ascolti.
    \end{itemize}
\end{itemize}

\paragraph{Autenticazione}
\begin{itemize}
    \item OAuth2 per integrazione Spotify/Apple Music.
    \item JWT per sessioni utente.
\end{itemize}

% -------------------------------------------------------------

\subsection{Frontend / Mobile}

La piattaforma deve offrire un’esperienza fluida su tutti i dispositivi, con un'interfaccia moderna e centrata sull’ascolto.

\paragraph{Opzioni consigliate}
\begin{itemize}
    \item \textbf{React Native} — ideale per app mobile multipiattaforma (iOS + Android).
    \item \textbf{React Web} — per dashboard, player web e Radio Ads Manager.
\end{itemize}

\paragraph{Motivazioni}
\begin{itemize}
    \item Codice condiviso mobile + web.
    \item Ecosistema solido di UI libraries (ex: Reanimated, Zustand/Redux).
    \item Ottima integrazione con WebSocket per Radio Rooms.
\end{itemize}

\paragraph{Componenti principali dell’interfaccia}
\begin{itemize}
    \item Player ibrido (TTS + Spotify/Apple).
    \item Editor di prompt.
    \item Visualizzazione Quote Cards.
    \item Radio Rooms (sincronizzazione + chat leggera opzionale).
    \item Modulo dedicato al Radio Ads Manager per advertiser.
\end{itemize}

% -------------------------------------------------------------

\subsection{AI \& Data}

Cuore logico e creativo della piattaforma, basato su modelli di linguaggio avanzati e pipeline di verifica.

\paragraph{LLM Provider}
\begin{itemize}
    \item \textbf{OpenAI GPT-5.1 / GPT-4.1} per:
    \begin{itemize}
        \item interpretazione prompt,
        \item generazione script,
        \item rubriche,
        \item creatività controllata.
    \end{itemize}
    \item Possibile uso di \textbf{modelli open-source} (Llama 3.x, Mistral) in modalità on-premise per ridurre costi a lungo termine.
\end{itemize}

\paragraph{TTS Provider}
\begin{itemize}
    \item \textbf{OpenAI TTS}, \textbf{ElevenLabs}, \textbf{Azure Neural TTS} per voce personalizzata e alta qualità.
    \item Pipeline duale \textit{Fast} + \textit{High Quality} per ottimizzare costi e latenza.
\end{itemize}

\paragraph{Vector DB}
\begin{itemize}
    \item \textbf{Pinecone / Weaviate / Qdrant} per:
    \begin{itemize}
        \item RAG,
        \item fact-checking,
        \item indexing delle curiosità verificate.
    \end{itemize}
\end{itemize}

\paragraph{Analytics \& Telemetry}
\begin{itemize}
    \item Mixpanel / Amplitude per analytics lato utente.
    \item Elastic / OpenSearch per logging distribuito.
    \item Prometheus + Grafana per monitoraggio e metriche tecniche.
\end{itemize}

% -------------------------------------------------------------

\subsection{Infrastructure}

La piattaforma necessita di un’infrastruttura elastica, affidabile e capace di gestire un elevato throughput audio.

\paragraph{Cloud Provider}
\begin{itemize}
    \item \textbf{AWS} (preferito) oppure \textbf{GCP}.
    \item Multi-zona per alta disponibilità.
\end{itemize}

\paragraph{Container Orchestration}
\begin{itemize}
    \item \textbf{Kubernetes (EKS / GKE)} per orchestrare i microservizi.
    \item \textbf{Helm} per il deployment.
\end{itemize}

\paragraph{Database}
\begin{itemize}
    \item \textbf{PostgreSQL} — utenti, episodi, campagne, profili.
    \item \textbf{Redis} — caching (prompt, intenti, playlist, chunk).
    \item \textbf{DynamoDB / MongoDB} — storage flexible per JSON (script e preferenze).
\end{itemize}

\paragraph{Storage}
\begin{itemize}
    \item \textbf{S3} per:
    \begin{itemize}
        \item file audio generati,
        \item asset utente (Quote Cards, profili),
        \item pre-rendering.
    \end{itemize}
\end{itemize}

\paragraph{CDN}
\begin{itemize}
    \item \textbf{CloudFront} per streaming veloce dei chunk vocali.
\end{itemize}

\paragraph{Queue / Event Bus}
\begin{itemize}
    \item \textbf{Kafka} o \textbf{RabbitMQ} per:
    \begin{itemize}
        \item orchestrazione chunk,
        \item gestione campagne Ads,
        \item aggiornamenti in tempo reale.
    \end{itemize}
\end{itemize}

\paragraph{Cache Layer}
\begin{itemize}
    \item Redis come caching ad alta velocità per ridurre latenza e costi dei modelli AI.
\end{itemize}

\paragraph{WebSocket Hub}
\begin{itemize}
    \item \textbf{AWS API Gateway WebSocket} oppure \textbf{Elixir/Phoenix Channels} per la sincronizzazione delle Radio Rooms.
\end{itemize}
% ---------------------------------
\section{Infrastructure, Costs \& Scalability}

Questa sezione fornisce una panoramica delle risorse infrastrutturali necessarie per il MVP e per la crescita graduale della piattaforma, insieme a una stima dei costi operativi e delle strategie adottate per garantire scalabilità, efficienza e affidabilità.

% -------------------------------------------------------------

\subsection{Baseline Infrastructure}

La seguente configurazione rappresenta un setup realistico per un MVP funzionante e in grado di supportare i primi utenti (10\,000--50\,000 sessioni mensili).

\paragraph{Compute}
\begin{itemize}
    \item 3--6 istanze \textbf{AWS EC2 t3.medium} / \textbf{t3.large} per i microservizi principali (API Gateway, Session Service, Playlist Service, Ads Manager).
    \item 2 istanze \textbf{c5.large} dedicate ai servizi AI interni (Prompt Interpreter, Fact-Checking, Script Generator).
    \item 1 istanza \textbf{g4dn.xlarge} opzionale per sperimentazioni di modelli open-source o per caching TTS avanzato.
\end{itemize}

\paragraph{Database \& Storage}
\begin{itemize}
    \item \textbf{PostgreSQL (RDS)} — entità core: utenti, episodi, playlist, campagne pubblicitarie.
    \item \textbf{Redis (ElastiCache)} — caching di prompt, JSON degli intenti, playlist, sessioni.
    \item \textbf{S3} — storage per:
        \begin{itemize}
            \item chunk vocali TTS,
            \item Quote Cards degli utenti,
            \item pre-rendering di intro/outro e template ricorrenti.
        \end{itemize}
\end{itemize}

\paragraph{Networking \& Distribution}
\begin{itemize}
    \item \textbf{CloudFront CDN} — distribuzione dei chunk vocali con bassa latenza.
    \item \textbf{API Gateway WebSocket} — sincronizzazione delle Radio Rooms.
    \item \textbf{Load Balancer (ALB)} per i microservizi.
\end{itemize}

\paragraph{Expected Usage (MVP)}
\begin{itemize}
    \item Episodi generati: 10\,000--30\,000/mese.
    \item Durata media episodio: 15--25 minuti.
    \item Dimensione media audio vocale per episodio: 8--20 MB.
    \item Richieste API: 3--5 milioni/mese.
\end{itemize}

% -------------------------------------------------------------

\subsection{AI Cost Estimation}

La generazione di un episodio comporta principalmente due costi:
\begin{enumerate}
    \item inferenze \textbf{LLM} per interpretazione prompt e generazione script;
    \item \textbf{TTS} per trasformare lo script in audio.
\end{enumerate}

\paragraph{1. LLM Costs}

Costo tipico per episodio (stimato):
\begin{itemize}
    \item Interpretazione prompt: 1--3 centesimi.
    \item Generazione script (10--20 chunk): 4--10 centesimi.
\end{itemize}

\textbf{Totale LLM per episodio: 5--13 centesimi}.  
Con ottimizzazioni (prompt compressi, caching): anche 3--7 centesimi.

\paragraph{2. TTS Costs}

Dipende dal provider (OpenAI, ElevenLabs, Azure TTS).  
Stima media:

\begin{itemize}
    \item 1 minuto di voce sintetica: 0.015--0.050 €.
\end{itemize}

Con un episodio medio di 3--8 minuti di parlato totale:
\[
\textbf{Costo TTS per episodio: 0.05--0.30 €}
\]

\paragraph{Costo totale per episodio (LLM + TTS)}

\[
\textbf{0.08--0.40 € per episodio}
\]

\paragraph{Scenari mensili}

\begin{itemize}
    \item 10\,000 episodi/mese → 800--4\,000 €
    \item 50\,000 episodi/mese → 4\,000--20\,000 €
    \item 200\,000 episodi/mese → 16\,000--80\,000 €
\end{itemize}

Riduzioni possibili:
\begin{itemize}
    \item caching di script ricorrenti e rubriche;
    \item uso di TTS ``Fast'' per ascolti non salvati;
    \item batch rendering notturno usando istanze spot.
\end{itemize}

% -------------------------------------------------------------

\subsection{Scaling Strategy}

La piattaforma adotta un insieme di meccaniche progettate per crescere a milioni di utenti attivi, mantenendo costi unitari stabili.

\paragraph{Chunked Rendering}
\begin{itemize}
    \item Gli script vengono generati in segmenti brevi (30--90s).
    \item Il TTS processa i chunk in pipeline: l'utente ascolta il primo chunk mentre i successivi vengono generati.
    \item Permette rigenerazioni veloci in caso di modifiche del prompt.
\end{itemize}

\paragraph{Caching dei contenuti AI}
\begin{itemize}
    \item Intro, rubriche o spiegazioni standard vengono pre-renderizzate in S3.
    \item Il Fact-Checking Engine memorizza curiosità già verificate.
    \item I chunk dei template (“Radio Notturna”, “Motivazionale”, ecc.) possono essere riutilizzati.
\end{itemize}

\paragraph{Autoscaling}
\begin{itemize}
    \item Kubernetes EKS scala automaticamente in base a CPU/RAM e code di messaggi.
    \item I servizi più costosi (TTS e Script Generator) possono essere assegnati a nodi dedicati.
\end{itemize}

\paragraph{CDN Distribution}
\begin{itemize}
    \item Tutti gli audio vocali vengono distribuiti globalmente tramite CloudFront.
    \item Riduce latenza, bandwidth server-side e costi di transito.
\end{itemize}

\paragraph{Prioritization \& Rate Control}
\begin{itemize}
    \item Parti “premium” o salvate hanno priorità maggiore rispetto all’anteprima rapida.
    \item I rendering HQ possono essere posticipati nelle ore di basso traffico.
\end{itemize}

% -------------------------------------------------------------

\subsection{Reliability \& Availability Targets}

L’applicazione segue principi SRE e mira agli standard di affidabilità tipici delle piattaforme audio.

\paragraph{SLO target}
\begin{itemize}
    \item \textbf{Availability:} 99.5\% (MVP) → 99.9\% (scalabilità).
    \item \textbf{Latency API:} $<$150 ms per richieste standard.
    \item \textbf{Start-to-play latency:} $<$2.0 s per l’avvio della puntata.
    \item \textbf{WebSocket Sync Drift:} $<$150 ms nelle Radio Rooms.
\end{itemize}

\paragraph{Backup \& Disaster Recovery}
\begin{itemize}
    \item Snapshot RDS giornalieri + replica multi-AZ.
    \item Backup incrementali S3.
    \item Configurazione “warm standby” per servizi critici (Session Service, Playlist Service).
\end{itemize}

\paragraph{Fault Tolerance}
\begin{itemize}
    \item Circuit breaker per i TTS provider (fallback automatico).
    \item Queue-based retry per generazioni chunk fallite.
    \item Health checks e restart automatico via Kubernetes.
\end{itemize}
% ---------------------------------
\section{API Design Overview}

Questa sezione descrive l’organizzazione generale delle API utilizzate nella piattaforma Radio-AI, distinguendo tra:
\begin{itemize}
    \item API esterne (public-facing), utilizzate dal client mobile/web;
    \item API interne tra microservizi, tipicamente gRPC o REST ad alte prestazioni;
    \item integrazioni con servizi terzi (musica, pagamenti, analytics).
\end{itemize}

L’obiettivo non è fornire la definizione completa degli endpoint, ma un quadro strutturale per comprendere il flusso dei dati e le responsabilità di ciascun modulo.

% -------------------------------------------------------------

\subsection{External APIs}

Le API esterne costituiscono il punto di ingresso dell’applicazione per gli utenti finali.  
Il design privilegia semplicità, sicurezza e basse latenze.

\paragraph{Autenticazione}
\begin{itemize}
    \item \verb|POST /auth/login| — login via provider esterno (Spotify, Apple, Google).
    \item \verb|POST /auth/refresh| — rinnovo token JWT.
    \item \verb|POST /auth/link-music-provider| — collega l’account Spotify/Apple per consentire l’uso delle playlist.
\end{itemize}

\paragraph{Generazione Episodi}
\begin{itemize}
    \item \verb|POST /episodes/create|  
    Payload:
    \begin{itemize}
        \item prompt dell’utente;
        \item playlist\_id (opzionale);
        \item parametri vocali;
        \item durata desiderata.
    \end{itemize}

    \item \verb|GET /episodes/{episode_id}|  
    Metadati dell’episodio (titolo, prompt rigenerato, durata, stato rendering).

    \item \verb|GET /episodes/{episode_id}/stream|  
    URL firmati per i chunk vocali e timeline musicale.

    \item \verb|PATCH /episodes/{episode_id}/prompt|  
    Aggiornamento live del prompt e rigenerazione dei chunk futuri.

    \item \verb|GET /episodes|  
    Lista episodi dell’utente (cronologia).
\end{itemize}

\paragraph{Radio Rooms}
\begin{itemize}
    \item \verb|POST /rooms/create| — crea una stanza sincronizzata.
    \item \verb|POST /rooms/{room_id}/join| — un utente entra nella stanza.
    \item \verb|GET /rooms/{room_id}/state| — stato corrente della sincronizzazione.
    \item \verb|WS /rooms/{room_id}/sync| — WebSocket per timestamp, comandi e drift correction.
\end{itemize}

\paragraph{Quote Cards \& Social Layer}
\begin{itemize}
    \item \verb|POST /quotes/create| — genera una Quote Card a partire da una frase della trascrizione.
    \item \verb|GET /quotes/{id}| — accede alla card.
    \item \verb|GET /profile/{user_id}| — informazioni pubbliche del profilo.
    \item \verb|GET /feed/public| — puntate e card condivise da altri utenti.
\end{itemize}

\paragraph{Ads Manager (per advertiser)}
\begin{itemize}
    \item \verb|POST /ads/campaigns| — crea una nuova campagna.
    \item \verb|GET /ads/campaigns/{id}| — dettagli campagna e andamento.
    \item \verb|PATCH /ads/campaigns/{id}| — aggiornamento targeting/budget.
    \item \verb|GET /ads/insights| — report aggregati.
\end{itemize}

% -------------------------------------------------------------

\subsection{Internal Service APIs}

I microservizi comunicano attraverso API interne ad alte prestazioni.  
La scelta consigliata è:
\begin{itemize}
    \item \textbf{gRPC} per comunicazioni sincrone ad alta efficienza (Script Generator, TTS Orchestrator, Playlist Service);
    \item \textbf{REST JSON} per servizi meno latenza-critical;
    \item \textbf{Event bus} (Kafka/RabbitMQ) per operazioni asincrone (rendering chunk, inserzioni, analytics).
\end{itemize}

\paragraph{Esempi di contratti ad alto livello}

\subsubsection*{Prompt Interpreter Service — gRPC}
\begin{itemize}
    \item \verb|InterpretPrompt(request PromptMessage) → IntentJSON|
\end{itemize}

\subsubsection*{Playlist Service — gRPC}
\begin{itemize}
    \item \verb|GetPlaylistMetadata(playlist_id)|
    \item \verb|GenerateAIPlaylist(intent_json)|
\end{itemize}

\subsubsection*{Script Generator Service — Event Driven}
Topic: \verb|script.generate|
\begin{itemize}
    \item Input: intent JSON + playlist JSON
    \item Output: lista chunk narrativi
\end{itemize}

\subsubsection*{TTS Orchestrator — gRPC/Queue}
\begin{itemize}
    \item \verb|SynthesizeChunk(chunk_text, voice_params)|
    \item Risultato: URL su S3 del file audio.
\end{itemize}

\subsubsection*{Streaming Orchestrator — REST}
\begin{itemize}
    \item \verb|GET /timeline/{episode_id}| — timeline di playback.
    \item \verb|POST /sync/update| — sincronizzazione timestamp.
\end{itemize}

\subsubsection*{Radio Rooms Service — WebSocket Hub}
\begin{itemize}
    \item Messaggi: \verb|PLAY|, \verb|PAUSE|, \verb|SEEK|, \verb|SYNC_UPDATE|.
\end{itemize}

\subsubsection*{Ads Targeting Engine — Event Bus}
Topic: \verb|ads.match.episode|
\begin{itemize}
    \item Input: metadati dell’episodio in generazione.
    \item Output: lista di inserzioni candidate.
\end{itemize}

% -------------------------------------------------------------

\subsection{Third-Party Integrations}

La piattaforma si appoggia a vari provider esterni, integrati tramite API sicure e tokenizzati.

\paragraph{Spotify / Apple Music}
\begin{itemize}
    \item Autenticazione OAuth.
    \item Accesso ai metadati playlist.
    \item Recommendation API (solo per Spotify).
    \item Riproduzione tramite SDK nativo del provider.
\end{itemize}

\paragraph{TTS Providers}
\begin{itemize}
    \item OpenAI TTS API.
    \item ElevenLabs API per voci premium.
    \item Azure Cognitive Services come backup provider.
\end{itemize}

\paragraph{LLM Provider}
\begin{itemize}
    \item OpenAI GPT (5.1 / 4.1) per:
    \begin{itemize}
        \item interpretazione prompt,
        \item generazione chunk narrativi,
        \item rubriche e curiosità culturali,
        \item riscrittura verificata durante fact-checking.
    \end{itemize}
\end{itemize}

\paragraph{Payment Provider}
\begin{itemize}
    \item Stripe — pagamenti per premium, abbonamenti, campagne Ads.
\end{itemize}

\paragraph{Auth Provider}
\begin{itemize}
    \item Firebase Auth, Auth0 oppure Cognito.
\end{itemize}

\paragraph{Analytics \& Crash Reporting}
\begin{itemize}
    \item Mixpanel / Amplitude — funnel e retention.
    \item Sentry — error monitoring mobile/web/backend.
    \item Google Analytics — traffico web.
\end{itemize}

\paragraph{CDN}
\begin{itemize}
    \item AWS CloudFront per distribuzione audio a bassa latenza.
\end{itemize}
% ---------------------------------
\section{Security \& Compliance}

La sicurezza e la conformità normativa rappresentano elementi centrali nella progettazione di Radio-AI.  
La piattaforma tratta dati sensibili (voci, preferenze musicali, prompt personali) e interagisce con provider esterni (Spotify, Apple Music, sistemi di pagamento), rendendo fondamentale un approccio robusto e verificabile alla protezione di utenti e contenuti.

% -------------------------------------------------------------

\subsection{Authentication \& Authorization}

La piattaforma implementa un modello di autenticazione moderno e multilivello, con sessioni sicure, token a scadenza controllata e integrazione con provider esterni.

\paragraph{Autenticazione}
\begin{itemize}
    \item \textbf{OAuth 2.0} per integrazione con Spotify, Apple, Google.
    \item \textbf{JWT (JSON Web Tokens)} per autenticazione verso le API interne.
    \item Token access e refresh con rotazione programmata.
    \item Pinning dei token ai dispositivi per ridurre il furto di sessione.
\end{itemize}

\paragraph{Autorizzazione}
\begin{itemize}
    \item Modello RBAC (Role-Based Access Control):
    \begin{itemize}
        \item ruolo utente standard (listener),
        \item ruolo creator (futuro),
        \item ruolo advertiser,
        \item ruolo admin.
    \end{itemize}
    \item Controlli granulari per:
    \begin{itemize}
        \item creazione e modifica episodi,
        \item gestione Radio Rooms,
        \item gestione campagne Ads.
    \end{itemize}
    \item Audit trail per tutte le operazioni sensibili.
\end{itemize}

\paragraph{Protezione delle credenziali musicali}
\begin{itemize}
    \item I token Spotify/Apple non vengono mai memorizzati in chiaro.
    \item KMS (Key Management Service) per cifratura AES256 lato server.
\end{itemize}

% -------------------------------------------------------------

\subsection{Privacy \& Data Protection}

Radio-AI rispetta la normativa europea GDPR e adotta un approccio di minimizzazione dei dati e trasparenza nel trattamento.

\paragraph{Tipologie di dati trattati}
\begin{itemize}
    \item Prompt testuali inseriti dall’utente.
    \item Trascrizioni generate dall’AI.
    \item File audio TTS generati (chunks vocali).
    \item Cronologia episodi ascoltati.
    \item Quote Cards e metadati social.
    \item Dati tecnici (ID playlist, timestamp, eventi player).
\end{itemize}

\paragraph{Data Minimization}
\begin{itemize}
    \item I prompt non vengono usati per addestrare modelli esterni.
    \item I chunk vocali sono archiviati solo se l’utente salva l’episodio.
    \item Tutti i dati non necessari al funzionamento vengono automaticamente rimossi dopo 30--90 giorni.
\end{itemize}

\paragraph{Data Retention \& Storage}
\begin{itemize}
    \item PostgreSQL per dati strutturati (cifrato at rest).
    \item S3 per file audio (server-side encryption con AES256).
    \item Log tecnici conservati 14--30 giorni, pseudonimizzati.
\end{itemize}

\paragraph{GDPR Compliance}
\begin{itemize}
    \item \textbf{Right to Access}: l’utente può esportare tutta la cronologia.
    \item \textbf{Right to Erasure}: cancellazione completa delle sessioni audio e dei prompt.
    \item \textbf{Data Portability}: esportazione episodi come JSON+metadata.
\end{itemize}

\paragraph{Sicurezza della Comunicazione}
\begin{itemize}
    \item HTTPS ovunque (TLS 1.2+).
    \item WebSocket sicuri (WSS) per le Radio Rooms.
    \item Strict-Transport-Security e protezione CSRF su tutte le form web.
\end{itemize}

% -------------------------------------------------------------

\subsection{Abuse Mitigation}

Per prevenire attività malevole, consumo anomalo di risorse e contenuti inappropriati, la piattaforma implementa livelli multipli di protezione.

\paragraph{Rate Limiting \& Throttling}
\begin{itemize}
    \item Limite sugli endpoint di creazione episodi.
    \item Riduzione priorità o blocco automatico quando vengono superate soglie di utilizzo sospette.
\end{itemize}

\paragraph{Prompt Moderation}
\begin{itemize}
    \item Sistema di moderazione AI per contenuti vietati (violenza esplicita, hate speech, illegalità).
    \item Flag dei prompt borderline per revisione manuale.
\end{itemize}

\paragraph{Protezione AI \& Anti-Automation}
\begin{itemize}
    \item ReCAPTCHA invisibile per prevenire bot.
    \item Analisi comportamentale per rilevare automazioni di massa sulla generazione episodi.
\end{itemize}

\paragraph{Abusi nelle Radio Rooms}
\begin{itemize}
    \item Kick e ban gestiti dall’host.
    \item Rate limit sui comandi WebSocket (play/pause/seek).
    \item Controllo drift e sincronizzazione handshake per evitare attacchi di flood.
\end{itemize}

\paragraph{Ads Manager Abuse Protection}
\begin{itemize}
    \item Validazione server-side degli script pubblicitari generati.
    \item Limiti di budget minimo e massimo per evitare spam.
    \item Verifica automatica del dominio/brand prima di accettare la campagna.
\end{itemize}
% ---------------------------------
\section{Roadmap}

La seguente roadmap suddivide l’evoluzione di Radio-AI in tre fasi principali:  
\textbf{MVP}, \textbf{Versione 1.0} e \textbf{sprint successivi post-lancio}.  
Gli obiettivi sono organizzati per priorità tecnica, impatto utente e sostenibilità economica.

% -------------------------------------------------------------

\subsection{MVP (0--6 Months)}

L’obiettivo dell’MVP è validare il core del prodotto: generazione di puntate radiofoniche personalizzate con musica proveniente da servizi esterni.

\paragraph{Obiettivi Chiave}
\begin{itemize}
    \item Generazione completa di una puntata con:
    \begin{itemize}
        \item interpretazione prompt;
        \item intent extraction;
        \item script generator;
        \item TTS rendering;
        \item sincronizzazione con playlist Spotify/Apple.
    \end{itemize}
    \item Player ibrido funzionante (musica esterna + voce AI).
    \item Supporto Modalità A (con playlist) e Modalità B (senza playlist).
    \item Gestione episodi:
    \begin{itemize}
        \item creazione,
        \item salvataggio,
        \item cronologia,
        \item rigenerazione parziale.
    \end{itemize}
    \item Prompt rigenerato (descrizione episodio).
    \item Sistema di autenticazione OAuth + JWT.
\end{itemize}

\paragraph{Milestones Tecniche}
\begin{enumerate}
    \item Setup infrastruttura cloud (K8s, DB, queue, S3).
    \item Integrazione Spotify/Apple Music.
    \item Implementazione API Gateway + microservizi core.
    \item TTS Orchestrator (Fast Mode).
    \item UX base dell’app mobile (React Native).
\end{enumerate}

\paragraph{Deliverable dell’MVP}
\begin{itemize}
    \item Un utente può generare e ascoltare un episodio completo.
    \item L’esperienza è stabile, fluida e ripetibile.
    \item Lato business: validazione iniziale con early testers.
\end{itemize}

% -------------------------------------------------------------

\subsection{Version 1.0 (6--12 Months)}

Questa fase introduce gli aspetti social, collaborativi e sperimentali, trasformando Radio-AI da semplice generatore in ecosistema di contenuti condivisi.

\paragraph{Nuove Funzionalità Utente}
\begin{itemize}
    \item \textbf{Radio Rooms}: ascolto sincronizzato multi-dispositivo/multi-utente.
    \item \textbf{Quote Cards}: selezione di frasi dalle puntate + generazione grafiche social-ready.
    \item \textbf{Profili Pubblici}: episodi condivisi, card, prompt salvati.
    \item \textbf{Feed Esplorazione}: ricerca per parola chiave, tema, mood, artista.
\end{itemize}

\paragraph{Espansioni Tecniche}
\begin{itemize}
    \item Fact-Checking Engine completo + caching avanzato.
    \item Miglioramento Script Generator (rubriche più complesse).
    \item TTS Orchestrator con \textbf{High-Quality Rendering} differito.
    \item Supporto multi-voce, parametri avanzati di personalizzazione.
\end{itemize}

\paragraph{Ads Manager (Beta)}
\begin{itemize}
    \item Creazione campagne.
    \item Targeting tematico.
    \item Generazione spot AI-based.
    \item Inserzioni automatiche per episodi coerenti.
\end{itemize}

\paragraph{Milestones Versione 1.0}
\begin{enumerate}
    \item Lancio soft pubblico (inviti).
    \item Integrazione analytics avanzati (retention, deep listening).
    \item Onboarding advertiser pilota.
\end{enumerate}

% -------------------------------------------------------------

\subsection{Beyond 1.0}

In questa fase Radio-AI evolve da piattaforma radio generata dall’AI a vero ecosistema narrativo, creativo e collaborativo.

\paragraph{Direzioni Future}
\begin{itemize}
    \item \textbf{Creator Studio}: strumenti per creare format personalizzati, template vocali, rubriche custom.
    \item \textbf{Live Streaming AI-assisted}: dirette “quasi live”, con generazione dinamica in tempo reale.
    \item \textbf{Marketplace di Voci Premium}: licensing di voci professionali (accordi con talent).
    \item \textbf{Partnership Editoriali}: etichette musicali, podcaster, radio tradizionali.
    \item \textbf{Distribuzione cross-platform}: esportazione come podcast, clip video TikTok-ready.
    \item \textbf{Collaborative Prompting}: puntate scritte a più mani dagli utenti.
    \item \textbf{Revenue Sharing per Creators}: monetizzazione tramite Ads Manager.
\end{itemize}

\paragraph{Obiettivi Strategici}
\begin{itemize}
    \item Stabilire Radio-AI come nuovo formato audio ibrido tra radio, podcast e AI.
    \item Scalare a milioni di utenti con costi per episodio ottimizzati.
    \item Creare un vantaggio competitivo sostenibile basato su:
    \begin{itemize}
        \item targeting tematico unico,
        \item contenuti personalizzati,
        \item ecosistema social proprietario.
    \end{itemize}
\end{itemize}
% ---------------------------------
\section{Open Questions \& Risks}

Nonostante la solidità della visione tecnica e del modello di prodotto, alcuni elementi cruciali richiedono ulteriori analisi, validazioni e decisioni strategiche.  
Questa sezione elenca le principali aree di incertezza, i rischi associati e i punti aperti da affrontare durante le fasi di sviluppo e scalabilità.

% -------------------------------------------------------------

\subsection{Licenze \& Aspetti Legali}

\paragraph{Musica}
\begin{itemize}
    \item L’utilizzo della musica avviene tramite \textbf{Spotify/Apple Music Connect} e non tramite riproduzione diretta della piattaforma.
    \item Rimangono aperte le seguenti questioni:
    \begin{itemize}
        \item stabilità nel lungo periodo delle API di terze parti (rischio di cambiamenti nei ToS);
        \item possibilità legale di sincronizzare parlato AI + musica proveniente da provider esterni;
        \item limiti territoriali o dispositivi supportati.
    \end{itemize}
\end{itemize}

\paragraph{Voci (TTS \& Personaggi Famosi)}
\begin{itemize}
    \item L’utilizzo di voci “ispirate” a personaggi noti può richiedere:
    \begin{itemize}
        \item diritto di immagine/voce,
        \item contratti con talent o eredi,
        \item limiti di utilizzo geografico e commerciale.
    \end{itemize}
    \item Rischio: richieste di rimozione o contestazioni se una voce è troppo simile senza licenza.
\end{itemize}

\paragraph{Contenuti generati}
\begin{itemize}
    \item Le rubriche devono essere basate su fatti verificati → rischio di errori o informazioni inesatte.
    \item Necessità di garantire che nessun contenuto violi copyright, marchi o proprietà intellettuale.
\end{itemize}

% -------------------------------------------------------------

\subsection{UX Uncertainty}

\paragraph{Accettazione del formato “radio AI”}
\begin{itemize}
    \item Non è chiaro come gli utenti reagiranno alla combinazione AI + musica personale.
    \item Possibile esigenza di iterare sul tono dello speaker, frequenza musicale, durata puntate.
\end{itemize}

\paragraph{Modifica live del prompt}
\begin{itemize}
    \item Potrebbe risultare complessa per utenti non tecnici.
    \item Necessità di un onboarding molto chiaro.
\end{itemize}

\paragraph{Radio Rooms}
\begin{itemize}
    \item Comportamenti sociali non prevedibili (troll, flood, ascolto disallineato).
    \item Sincronizzazione perfetta tra dispositivi diversi potrebbe risultare difficile da percepire o valutare dall’utente.
\end{itemize}

% -------------------------------------------------------------

\subsection{Technical Unknowns}

\paragraph{Carico AI in scenari di crescita}
\begin{itemize}
    \item I costi TTS possono crescere rapidamente oltre una certa soglia di utenti attivi.
    \item Rischio di saturazione del rendering se la domanda supera il throughput.
\end{itemize}

\paragraph{Sincronizzazione audio}
\begin{itemize}
    \item Playback ibrido (musica locale + voce AI remota) richiede tolleranze di latenza molto strette.
    \item Differenze tra dispositivi, reti e provider musicali possono introdurre drift difficili da compensare.
\end{itemize}

\paragraph{Caching avanzato}
\begin{itemize}
    \item Non è scontato quanto contenuto possa essere riutilizzato tra utenti diversi.
    \item Correlazione tra temi, rubriche e prompt potrebbe essere inferiore al previsto.
\end{itemize}

% -------------------------------------------------------------

\subsection{Business \& Monetization Risks}

\paragraph{Ads Manager}
\begin{itemize}
    \item L’efficacia del targeting tematico deve essere validata con advertiser reali.
    \item Rischio di domanda insufficiente nelle prime fasi.
    \item Rischio regolatorio legato alla pubblicità audio personalizzata.
\end{itemize}

\paragraph{Costo unitario per episodio}
\begin{itemize}
    \item Potrebbe risultare troppo alto per una monetizzazione tradizionale (ads o subscription).
    \item Necessità di ottimizzare TTS e LLM per garantire margini sostenibili.
\end{itemize}

% -------------------------------------------------------------

\subsection{Compliance \& Regulatory Risks}

\paragraph{GDPR}
\begin{itemize}
    \item La gestione di prompt personali e voci potrebbe essere considerata sensibile.
    \item Richieste di cancellazione dei dati devono essere gestite con estrema precisione.
\end{itemize}

\paragraph{Contenuti AI}
\begin{itemize}
    \item Normative future sull’AI generativa potrebbero richiedere:
    \begin{itemize}
        \item watermarking dei contenuti vocali,
        \item segnalazione esplicita dell’uso di AI,
        \item limiti sull’uso di voci realistiche.
    \end{itemize}
\end{itemize}

% -------------------------------------------------------------

\subsection{Operational Risks}

\paragraph{Dipendenza da provider esterni}
\begin{itemize}
    \item Spotify/Apple possono cambiare API, costi o limitazioni.
    \item TTS provider possono aumentare i prezzi o limitare volumi massimi.
\end{itemize}

\paragraph{Failure di componenti critici}
\begin{itemize}
    \item Il TTS Orchestrator è un single point of failure se non ridondato.
    \item Il WebSocket Hub deve garantire alta disponibilità per Radio Rooms.
\end{itemize}

% -------------------------------------------------------------

\subsection{Mitigazione (alto livello)}

\begin{itemize}
    \item Contratti e SLA con provider AI e musicali.
    \item Caching aggressivo di contenuti vocali ricorrenti.
    \item Circuit breaker e fallback su provider TTS alternativi.
    \item Interfacce modulari per cambiare modello LLM rapidamente.
    \item Team legale per gestione licensing e diritti vocali.
    \item Iterazione UX continua tramite A/B testing.
\end{itemize}
% ---------------------------------
\section{Conclusion}

Radio-AI introduce un nuovo paradigma nell’esperienza audio: un ecosistema capace di combinare musica personale, narrazione generata dall’intelligenza artificiale, contenuti verificati e interazioni sociali in un formato che non esisteva prima.  
La piattaforma permette agli utenti di trasformare qualsiasi playlist in una diretta radiofonica personalizzata, dinamica e profondamente adattata al contesto di ascolto, aprendo allo stesso tempo nuove opportunità creative e commerciali.

Il valore chiave di Radio-AI risiede nella sua architettura modulare e scalabile:
\begin{itemize}
    \item una pipeline AI avanzata per interpretazione prompt, generazione narrativa e sintesi vocale;
    \item un’infrastruttura progettata per sostenere milioni di sessioni con costi unitari ottimizzabili;
    \item strumenti sociali e collaborativi che trasformano episodi individuali in contenuti condivisibili;
    \item un sistema pubblicitario guidato dal targeting tematico, capace di valorizzare il contesto emotivo dell’ascolto.
\end{itemize}

Guardando al futuro, la roadmap prevede l’espansione verso funzionalità più profonde: un Creator Studio per format avanzati, un marketplace di voci premium, partnership editoriali, e forme di streaming “quasi-live” assistite dall’AI.  
Questi sviluppi, uniti alla struttura tecnica già delineata, posizionano Radio-AI come una piattaforma di riferimento per la prossima generazione di contenuti audio personalizzati.

Il prossimo passo è tradurre la visione in un prodotto concreto: un MVP solido, iterazioni rapide basate sul feedback reale e una strategia di crescita mirata che valorizzi l’unicità della proposta.  
Radio-AI non è solamente un’app: è un nuovo modo di vivere l’audio, modellato dall’utente, alimentato dall’AI e potenzialmente scalabile a livello globale.

\section{Appendix A: Example Prompts}
\label{appendix:prompts}

Questa sezione raccoglie alcuni esempi di prompt reali per illustrare la varietà di scenari supportati da Radio-AI. Per ciascun esempio riportiamo il testo del prompt (in forma sintetica) e gli intent principali che il sistema è chiamato a riconoscere.

\subsection*{A.1 Motivazione personale prima di un appuntamento}

\textbf{Prompt (esempio):}
\begin{quote}
Ciao, sto andando ad un appuntamento e sono molto nervoso. Crea una diretta radiofonica per darmi la carica e aiutarmi a sentirmi più sicuro. Includi telefonate di persone che raccontano la loro esperienza con i primi appuntamenti per confortarmi. Usa principalmente queste canzoni [lista] e, ogni tanto, inserisci qualche breve curiosità sulle indiscrezioni riguardo la nuova PS6.
\end{quote}

\textbf{Intent principali:}
\begin{itemize}[leftmargin=*]
    \item \textbf{Tema}: primo appuntamento, motivazione personale.
    \item \textbf{Rubriche}: telefonate immaginarie di ascoltatori con esperienze analoghe.
    \item \textbf{Curiosità reali}: news brevi e verificate su PS6.
    \item \textbf{Musica}: playlist guidata dall’utente.
    \item \textbf{Tono speaker}: energico, incoraggiante, empatico.
\end{itemize}

\subsection*{A.2 Supporto genitoriale per inizio scuola}

\textbf{Prompt (esempio):}
\begin{quote}
Ciao, mia figlia sta per iniziare le elementari e io vorrei esserle di supporto. Crea una diretta che affronti questo tema, includendo una conversazione con una psicologa che parli sia dal punto di vista del genitore che del bambino.
\end{quote}

\textbf{Intent principali:}
\begin{itemize}[leftmargin=*]
    \item \textbf{Tema}: supporto emotivo genitoriale, inizio scuola.
    \item \textbf{Rubriche}: finta intervista con una psicologa (contenuto basato su linee guida reali).
    \item \textbf{Musica}: selezione a discrezione dell’AI, coerente con mood rassicurante.
    \item \textbf{Tono speaker}: calmo, empatico, rassicurante.
\end{itemize}

\subsection*{A.3 Ripasso accademico con musica minima}

\textbf{Prompt (esempio):}
\begin{quote}
Tra un paio d’ore ho un esame di Analisi I. Genera una diretta per aiutarmi a ripassare, con linguaggio tecnico-scientifico. Metti poca musica: ogni venti minuti di spiegazione metti una canzone di David Bowie. Durante il parlato, tieni un sottofondo di musica classica, preferibilmente Chopin.
\end{quote}

\textbf{Intent principali:}
\begin{itemize}[leftmargin=*]
    \item \textbf{Tema}: ripasso accademico (Analisi I).
    \item \textbf{Rubriche}: spiegazioni teoriche, esempi d’esercizio, mini-riassunti.
    \item \textbf{Musica}: Bowie ogni 20 minuti, sottofondo classico continuo.
    \item \textbf{Tono speaker}: tecnico, chiaro, orientato allo studio.
\end{itemize}

\subsection*{A.4 Star Wars vs Star Trek}

\textbf{Prompt (esempio):}
\begin{quote}
Genera una diretta sulla sfida ``Meglio: Star Wars o Star Trek?''. Nel corso della puntata prova ad eleggere la saga migliore sulla base di fattori cinematografici e socio-culturali. Inserisci le colonne sonore più iconiche, presentandole con la storia della composizione e qualche curiosità.
\end{quote}

\textbf{Intent principali:}
\begin{itemize}[leftmargin=*]
    \item \textbf{Tema}: confronto tra saghe, analisi cinematografica.
    \item \textbf{Rubriche}: segmenti di approfondimento storico-culturale, analisi dei personaggi, impatto sul fandom.
    \item \textbf{Curiosità reali}: aneddoti verificati sulle colonne sonore e sui compositori.
    \item \textbf{Musica}: soundtrack selezionate da Star Wars e Star Trek.
    \item \textbf{Tono speaker}: narrativo, leggermente giocoso, ma informato.
\end{itemize}

\subsection*{A.5 Epopea fantasy personalizzata}

\textbf{Prompt (esempio):}
\begin{quote}
Genera una diretta che narri le gesta eroiche (con linguaggio medievale) di Matteo ``La Foglia'' e dei suoi compagni Linzo, Taxi e Wiz, che devono sconfiggere un drago che minaccia il villaggio di Castiglione. Dividi il racconto in atti, ognuno seguito da una canzone da locanda che richiami ciò che è appena successo, includendo ``La guerra di Piero'' di De André e ``Drunken Sailor''.
\end{quote}

\textbf{Intent principali:}
\begin{itemize}[leftmargin=*]
    \item \textbf{Tema}: racconto fantasy personalizzato.
    \item \textbf{Rubriche}: atti narrativi + dialoghi tra personaggi (fiction pura).
    \item \textbf{Musica}: brani medievali/folk + canzoni richieste.
    \item \textbf{Tono speaker}: storyteller, epico, ironico.
\end{itemize}

\subsection*{A.6 Viaggio verso l’aeroporto – Speciale Fuerteventura}

\textbf{Prompt (esempio):}
\begin{quote}
Sono in autostrada con gli amici per un viaggio di 1h20 verso l’aeroporto, volo per Fuerteventura. Genera una diretta di quella durata con consigli su luoghi da visitare, paesaggi, fauna, storia, personaggi famosi, leggende e qualche curiosità politica. Usa una playlist Spotify selezionata da me: ogni 1--3 minuti di parlato metti una canzone, e ogni 25 minuti fai un blocco di almeno 3 canzoni di fila.
\end{quote}

\textbf{Intent principali:}
\begin{itemize}[leftmargin=*]
    \item \textbf{Tema}: viaggio, guida culturale su Fuerteventura.
    \item \textbf{Curiosità reali}: informazioni verificate su storia, fauna, geografia, cultura locale.
    \item \textbf{Musica}: playlist utente, con pattern di alternanza parlato/musica preciso.
    \item \textbf{Durata}: puntata time-boxed (80 minuti).
\end{itemize}

\subsection*{A.7 Shuffle totale + citazioni da film}

\textbf{Prompt (esempio):}
\begin{quote}
Seleziono una playlist da Spotify e voglio una diretta senza speaker: solo musica in shuffle e citazioni iconiche tratte da film prima di ogni brano. Ogni 8 canzoni consigliami una canzone aggiuntiva simile alle precedenti, non presente in playlist.
\end{quote}

\textbf{Intent principali:}
\begin{itemize}[leftmargin=*]
    \item \textbf{Tema}: esperienza musicale quasi-pura, con layer cinematografico.
    \item \textbf{Rubriche}: citazioni da film (contenuto testuale breve).
    \item \textbf{Musica}: playlist utente in shuffle + raccomandazioni AI-based.
    \item \textbf{Tono speaker}: assente (solo voice-over delle citazioni, se richiesto).
\end{itemize}

\section{Appendix B: Terminology}
\label{appendix:terminology}

Questa appendice definisce alcuni termini chiave utilizzati all’interno del whitepaper.

\begin{description}[leftmargin=!,labelwidth=3cm]
    \item[Episodio / Episode] Una singola diretta radiofonica generata da Radio-AI a partire da un prompt e/o da una playlist. Include script, chunk audio e metadati (tema, mood, durata, ecc.).
    
    \item[Prompt] Testo inserito dall’utente nella barra di input (breve o lungo, strutturato o colloquiale) che descrive tema, preferenze, richieste e vincoli per l’episodio.
    
    \item[Prompt Normalizzato] Versione riscritta e strutturata del prompt originale, generata dall’AI per fungere da descrizione ufficiale dell’episodio (usata per catalogazione, ricerca, condivisione).
    
    \item[Intent] Insieme di istruzioni semantiche estratte dal prompt: tema principale, rubriche, mood e stile dello speaker, preferenze musicali, vincoli di durata, ecc.
    
    \item[Rubrica] Segmento narrativo ricorrente all’interno di una puntata (es. curiosità sugli artisti, mini-documentari, telefonate di ascoltatori immaginari, what-if, interviste simulate).
    
    \item[Chunk] Unità atomica di contenuto parlato all’interno di un episodio (intro, bridge, rubrica, ad, outro, ecc.), con script e relativo file audio TTS. I chunk sono generati e riprodotti in sequenza.
    
    \item[Playlist] Insieme di brani musicali utilizzato come base per la scaletta dell’episodio, proveniente dal servizio esterno (es. Spotify, Apple Music) o generato dall’AI.
    
    \item[Room / Radio Room] Sessione di ascolto condivisa in cui più utenti ascoltano lo stesso episodio in modo sincronizzato. Una Room è associata a un episodio e ha un host che ne controlla le impostazioni.
    
    \item[Quote Card] Asset grafico generato a partire da un estratto testuale della puntata (trascrizione). È pensato per la condivisione social e può includere testo, colori, font e altri elementi visivi.
    
    \item[Profilo Utente] Vista aggregata delle attività di un utente: episodi generati, prompt salvati, Quote Cards, impostazioni di privacy, relazioni social (follower/following).
    
    \item[Campagna / Campaign] Insieme di parametri pubblicitari creati da un advertiser in Radio Ads Manager (obiettivo, budget, durata, targeting tematico, creatività audio).
    
    \item[Spot / Ad Spot] Singola inserzione audio associata a una campagna, inserita in punti specifici degli episodi (intro, midroll, pre/post brano, outro) e misurata tramite impression e completion rate.
    
    \item[Targeting Tematico] Modello di targeting pubblicitario basato su temi, mood, generi musicali e atmosfere degli episodi, in contrasto con il targeting puramente demografico.
    
    \item[TTS Orchestrator] Componente che coordina le richieste di sintesi vocale verso uno o più provider TTS, gestendo caching, quality tier e fallback in caso di errore.
    
    \item[Fact-Checking Engine] Pipeline dedicata a verificare la correttezza di curiosità, dati storici e riferimenti culturali tramite modelli LLM e fonti affidabili; filtra o riscrive contenuti non verificabili.
    
    \item[LLM] Large Language Model utilizzato per interpretare i prompt, generare script, rubriche, dialoghi e descrizioni strutturate degli episodi.
    
    \item[Room Host] Utente che crea una Radio Room e controlla diritti di modifica del prompt, permessi per altri partecipanti e stato della stanza (attiva/chiusa).
\end{description}

\section{Appendix C: References}
\label{appendix:references}

Questa sezione elenca alcune categorie di riferimenti utili per l’implementazione di Radio-AI. I riferimenti specifici (URL, versioni delle API, paper tecnici) andranno aggiornati in base alle scelte effettive di stack tecnologico.

\subsection*{C.1 API e Servizi Esterni}

\begin{itemize}[leftmargin=*]
    \item Documentazione ufficiale dei provider LLM utilizzati (ad es. API di modelli generativi per prompt interpretation, script generation e supporto RAG).
    \item Documentazione dei provider TTS neurali (synthetic voice, voice cloning, parametri di timbro e prosodia).
    \item Documentazione Spotify Web API e/o Apple Music API per:
    \begin{itemize}
        \item accesso alle playlist utente;
        \item riproduzione via Connect/SDK;
        \item gestione delle autorizzazioni OAuth.
    \end{itemize}
    \item Documentazione del provider di pagamento (es. Stripe, Adyen, ecc.) per gestione subscription, crediti e fatturazione delle campagne Ads.
    \item Documentazione del provider di analytics (es. Mixpanel, Amplitude, Segment) per eventi di ascolto, retention, funnel.
\end{itemize}

\subsection*{C.2 Infrastruttura \& Architetture Cloud}

\begin{itemize}[leftmargin=*]
    \item Linee guida del cloud provider scelto (es. AWS, GCP, Azure) su:
    \begin{itemize}
        \item microservizi containerizzati (ECS/Kubernetes/GKE);
        \item storage oggetti per file audio (S3/Cloud Storage);
        \item CDN per distribuzione globale degli asset audio;
        \item sistemi di messaggistica/queue (SQS, Pub/Sub, Kafka).
    \end{itemize}
    \item Best practice per architetture event-driven e streaming audio a bassa latenza (WebSocket, WebRTC, HLS a segmenti corti).
\end{itemize}

\subsection*{C.3 LLM, TTS e Fact-Checking}

\begin{itemize}[leftmargin=*]
    \item Paper e documentazione su Large Language Models applicati a:
    \begin{itemize}
        \item few-shot prompt engineering;
        \item intent detection;
        \item narrative generation.
    \end{itemize}
    \item Letteratura su Text-to-Speech neurale (Tacotron, VITS, modelli neural codec) e controlli prosodici.
    \item Materiale su Retrieval-Augmented Generation (RAG) per fact-checking:
    \begin{itemize}
        \item vector search;
        \item ranking di fonti;
        \item mitigazione di allucinazioni.
    \end{itemize}
\end{itemize}

\subsection*{C.4 Legal, Privacy \& Compliance}

\begin{itemize}[leftmargin=*]
    \item Linee guida GDPR per trattamento di dati personali, log di ascolto, cronologia, preferenze.
    \item Documentazione su gestione consensi per:
    \begin{itemize}
        \item utilizzo di voci personalizzate;
        \item eventuale user-generated voice cloning;
        \item condivisione pubblica di episodi e Quote Cards.
    \end{itemize}
    \item Materiale di riferimento su licensing musicale (SIAE, Soundreef, equivalenti esteri) e modelli di integrazione via account dell’utente (es. analoghi a Spotify Connect).
\end{itemize}

\end{document}